\documentclass[11pt]{article}

\usepackage[margin=1in]{geometry}
\usepackage{amsmath,amssymb,amsthm,mathtools}
\usepackage{enumitem}
\usepackage{booktabs}
\usepackage{listings}
\usepackage[hidelinks]{hyperref}

\newtheorem{theorem}{Theorem}
\newtheorem{lemma}[theorem]{Lemma}
\newtheorem{proposition}[theorem]{Proposition}

\theoremstyle{definition}
\newtheorem{definition}[theorem]{Definition}
\newtheorem{assumption}[theorem]{Assumption}

\theoremstyle{remark}
\newtheorem{remark}[theorem]{Remark}

\lstset{
  basicstyle=\ttfamily\small,
  breaklines=true,
  frame=single,
  columns=fullflexible
}

\lstdefinelanguage{TLA}{
  morekeywords={VARIABLES,VARIABLE,LET,IN,IF,THEN,ELSE,UNCHANGED,
    EXTENDS,CONSTANT,CONSTANTS,ASSUME,THEOREM,Next,Init,Spec},
  morecomment=[l]{\\*},
  sensitive=true,
  morestring=[b]",
}

\lstdefinelanguage{proto}{
  morekeywords={syntax,package,message,enum,bytes,uint32,uint64,bool,repeated,reserved},
  morecomment=[l]{//},
  sensitive=true,
  morestring=[b]",
}

\newcommand{\concat}{\ensuremath{\mathbin{\Vert}}}
\newcommand{\bits}{\ensuremath{\{0,1\}}}
\newcommand{\Hf}{\ensuremath{\mathsf{H}}}
\newcommand{\HKDF}{\ensuremath{\mathsf{HKDF}}}
\newcommand{\MACD}{\ensuremath{\mathsf{MACANDD}}}
\newcommand{\StepKeyGen}{\ensuremath{\mathsf{StepKeyGen}}}
\newcommand{\StepSign}{\ensuremath{\mathsf{StepSign}}}
\newcommand{\StepVerify}{\ensuremath{\mathsf{StepVerify}}}
\newcommand{\PartSign}{\ensuremath{\mathsf{PartSign}}}
\newcommand{\PartVerify}{\ensuremath{\mathsf{PartVerify}}}
\newcommand{\SigCommit}{\ensuremath{\mathsf{SigCommit}}}
\newcommand{\ctx}[1]{\texttt{"#1"}}
\newcommand{\WOTS}{\textsc{wots}}

\newcommand{\Tropic}{\textsc{TROPIC01}}
\newcommand{\Pico}{\textsc{Raspberry Pi Pico 2 W}}
\newcommand{\RP}{\textsc{RP2350}}
\newcommand{\Lthree}{\ensuremath{\mathsf{L3}}}

\newcommand{\Cert}{\ensuremath{\mathsf{Cert}}}
\newcommand{\Pkg}{\ensuremath{\mathsf{Pkg}}}
\newcommand{\Active}{\ensuremath{\mathsf{Active}}}
\newcommand{\Ready}{\ensuremath{\mathsf{Ready}}}
\newcommand{\Prepared}{\ensuremath{\mathsf{Prepared}}}
\newcommand{\Committed}{\ensuremath{\mathsf{Committed}}}
\newcommand{\Accept}{\ensuremath{\mathsf{Accept}}}
\newcommand{\Reject}{\ensuremath{\mathsf{Reject}}}
\newcommand{\Online}{\ensuremath{\mathsf{Online}}}

\newcommand{\Bundle}{\ensuremath{B}}
\newcommand{\Ahead}{\ensuremath{A}}
\newcommand{\Bhead}{\ensuremath{J}}
\newcommand{\BootTicket}{\ensuremath{\mathsf{BootTicket}}}
\newcommand{\BootChain}{\ensuremath{\mathsf{BootChain}}}
\newcommand{\pkhw}{\ensuremath{\mathsf{pk}_{\mathrm{hw}}}}
\newcommand{\skhw}{\ensuremath{\mathsf{sk}_{\mathrm{hw}}}}
\newcommand{\pkh}{\ensuremath{P_{\mathrm{hw}}}}

\newcommand{\MACDcmd}{\texttt{MAC\_And\_Destroy}}
\newcommand{\InitCmd}{\texttt{MCounter\_Init}}
\newcommand{\GetCmd}{\texttt{MCounter\_Get}}
\newcommand{\UpdateCmd}{\texttt{MCounter\_Update}}
\newcommand{\Handle}{\ensuremath{\mathsf{handle}}}

\title{\textbf{Boot Fenced Fused Anchor Authority\\
for DSM Offline Bearer State}\\[4pt]
\large One Way Birth Fuse, Fused Anchor Head, RP2350 Partition Witness,\\
TROPIC01 MACANDD and Counter Evidence, Receiver Challenge,\\
Public DSM Verification, Recovery, and Tripwire Reconciliation\\
on Raspberry Pi Pico 2 W}
\author{Brandon Ramsay (Cryptskii)\\\small Irrefutable Labs Inc.}
\date{July 2026}

\begin{document}

\maketitle

\begin{abstract}
A Deterministic State Machine, or DSM, advances state by local deterministic acceptance rather
than by global consensus. Ordinary DSM operation does not require a blockchain, validator set,
sequencer, wall clock, or online settlement step on the common path.

Offline bearer transfer is the hard case. A receiver accepts a transfer while offline. The receiver
must not accept copied software pretending to be the live appliance, must not accept a second use
of the same spendable parent, and must not allow copied enrollment data to resume on new hardware.

This paper specifies a compact offline bearer authority for DSM using a Raspberry Pi Pico 2 W,
the RP2350 secure partition, and a MIKROE 6559 Secure Tropic Click carrying a TROPIC01 secure
element. The design does not put all trust in TROPIC01 and does not put all trust in the RP2350
partition. Instead, DSM, the partition, and TROPIC01 are fused into one forward only lineage.

The authority is based on four rules:

\[
\text{destroy the birth preimage,}
\]

\[
\text{boot fence the appliance before offline use,}
\]

\[
\text{advance one DSM root and one fused anchor head together,}
\]

\[
\text{accept only authenticated TROPIC01 counter evidence and public DSM validity.}
\]

The DSM root commits to an immutable anchor bundle \(\Bundle\), a fused anchor head \(\Ahead_i\),
a boot head \(\Bhead_b\), and an anchor counter \(u_i\). The anchor bundle binds the RP2350
partition key, TROPIC01 anchor identifier, enrolled counter value \(H_0\), MACANDD slots, device
identifier, policy hash, and the hash of a destroyed one way birth fuse. The fused anchor head
binds the DSM root lineage, the partition lineage, and the TROPIC01 MACANDD and counter lineage.

Every boot must produce a boot ticket chained from the DSM committed boot head. Every offline
release must bind the current boot ticket. A copied state image on new hardware cannot resume
offline bearer mode because new hardware cannot advance the committed boot head under the enrolled
anchor bundle.

A release is accepted only if the receiver verifies the DSM transition, the receiver challenge,
the RP2350 partition certificate, the TROPIC01 MACANDD witness, authenticated TROPIC01 counter
evidence, and the fused anchor head update. TROPIC01 evidence is necessary but not sufficient.
RP2350 partition evidence is necessary but not sufficient. DSM validity remains public and
receiver verified.
\end{abstract}

\section{Purpose and Scope}

This document specifies an optional DSM offline bearer authority. It is used only when a receiver
accepts a DSM transfer without online reconciliation at the moment of exchange.

The authority provides:

\begin{enumerate}[label=(\arabic*),leftmargin=2em]
\item one active SMT root for the appliance;
\item one immutable anchor bundle;
\item one forward only fused anchor head;
\item one boot fenced offline session head;
\item one physical counter derived anchor counter;
\item one certified root advance at a time;
\item receiver challenge binding;
\item public verification of the DSM transition;
\item RP2350 secure partition release evidence;
\item TROPIC01 hardware presence through MACANDD;
\item TROPIC01 counter evidence not trusted through the host;
\item recovery by re emitting the same committed root advance;
\item online fallback when local state, counter evidence, boot evidence, or policy does not match;
\item Tripwire exposure of any fork on reconciliation.
\end{enumerate}

The target hardware is:

\begin{center}
\begin{tabular}{lll}
\toprule
\textbf{Layer} & \textbf{Part} & \textbf{Role}\\
\midrule
Controller & Raspberry Pi Pico 2 W & host and transport board\\
MCU & RP2350 & secure partition and appliance policy\\
Secure element board & MIKROE 6559 Secure Tropic Click & TROPIC01 over SPI\\
Secure element & TROPIC01 & MACANDD, counter, R memory, pairing policy\\
Interface & SPI at 3.3 V & secure element command transport\\
\bottomrule
\end{tabular}
\end{center}

The design does not require the RP2350 to verify the entire DSM transition. The receiver verifies
DSM validity. The RP2350 partition stores appliance state, signs partition certificates, advances
the boot ratchet, and drives TROPIC01. The receiver does not accept a root advance merely because
the RP2350 says it happened.

The design also does not let TROPIC01 authorize DSM state. TROPIC01 contributes hardware witness
output and authenticated counter evidence. DSM validity remains public and receiver verified.

\section{Design Summary}

The appliance state includes:

\[
\Active =
(
h_i,
\Bundle,
\Ahead_i,
\Bhead_b,
u_i,
\mathsf{status},
\mathsf{record}
).
\]

Here:

\begin{itemize}
\item \(h_i\) is the active DSM SMT root;
\item \(\Bundle\) is the immutable anchor bundle;
\item \(\Ahead_i\) is the active fused anchor head;
\item \(\Bhead_b\) is the last DSM committed boot head;
\item \(u_i\) is the active anchor counter;
\item \(\mathsf{status}\in\{\Ready,\Prepared,\Committed\}\);
\item \(\mathsf{record}\) is empty, prepared, or committed.
\end{itemize}

The previous DSM state commits to:

\[
(
\Bundle,
\Ahead_i,
\Bhead_b,
u_i
).
\]

A valid offline transfer from \(h_i\) must advance to a next root \(h_{i+1}\) that commits to:

\[
(
\Bundle,
\Ahead_{i+1},
\Bhead_{b'},
u_i+1
).
\]

The release is not merely:

\[
(h_i,u_i)\rightarrow(h_{i+1},u_i+1).
\]

The fused form is:

\[
(h_i,\Ahead_i,\Bhead_b,u_i)
\rightarrow
(h_{i+1},\Ahead_{i+1},\Bhead_{b'},u_i+1).
\]

The receiver accepts only if:

\begin{enumerate}[label=(\arabic*),leftmargin=2em]
\item the previous root is the receiver accepted root for the object, under the accepted root
frontier rule (adoption at relationship genesis, then strict continuity);
\item the previous state commits to \(\Bundle,\Ahead_i,\Bhead_b,u_i\);
\item a boot ticket or boot chain advances \(\Bhead_b\) to the current boot head \(\Bhead_{b'}\);
\item the claimed next anchor counter is \(u_i+1\);
\item the DSM proof verifies \(h_i\rightarrow h_{i+1}\);
\item the next root commits to \(\Bundle,\Ahead_{i+1},\Bhead_{b'},u_i+1\);
\item the RP2350 partition certificate verifies over the same root advance message;
\item the TROPIC01 MACANDD witness verifies over the same root advance message and partition
commitment;
\item authenticated TROPIC01 counter evidence shows the live physical counter corresponds to
\(u_i+1\);
\item the receiver challenge matches;
\item no policy event invalidates the anchor.
\end{enumerate}

The core rule is:

\[
\boxed{
\text{The receiver accepts only when DSM, the RP2350 partition, and TROPIC01 all bind the same
root advance and fused anchor head.}
}
\]

\section{Naming Discipline}

The protocol uses the following names.

\begin{center}
\begin{tabular}{ll}
\toprule
\textbf{Name} & \textbf{Meaning}\\
\midrule
\texttt{prev\_root} & SMT root advanced from\\
\texttt{next\_root} & SMT root advanced to\\
\texttt{anchor\_counter} & derived increasing DSM value \(u=H_0-H\)\\
\texttt{next\_anchor\_counter} & required next value \(u+1\)\\
\texttt{enrolled\_counter} & enrolled TROPIC01 physical counter value \(H_0\)\\
\texttt{attested\_live\_counter} & raw authenticated TROPIC01 value \(H\)\\
\texttt{anchor\_bundle} & immutable enrollment digest \(\Bundle\)\\
\texttt{anchor\_head} & fused anchor head \(\Ahead_i\)\\
\texttt{boot\_head} & boot fence head \(\Bhead_b\)\\
\bottomrule
\end{tabular}
\end{center}

The anchor counter is not the raw TROPIC01 counter. The anchor counter increases:

\[
u=H_0-H.
\]

The raw TROPIC01 counter counts down:

\[
H\leftarrow H-1.
\]

The receiver computes:

\[
u_{\mathrm{attested}}=H_0-H_{\mathrm{attested}}.
\]

The receiver accepts only if:

\[
u_{\mathrm{attested}}=u_i+1.
\]

\section{Why a Counter Alone Is Not the Design}

A chip signature over transfer details and a counter gives an ordered hardware log. It does not
by itself prove that a proposed transfer is a valid DSM successor from the receiver accepted
previous root. It also does not prove that the release came from the correct fused appliance
lineage.

The useful object is not:

\[
\mathsf{Sign}_k(\mathsf{transfer}\concat u).
\]

The useful object is:

\[
(h_i,\Ahead_i,\Bhead_b,u_i)
\rightarrow
(h_{i+1},\Ahead_{i+1},\Bhead_{b'},u_i+1),
\]

where \(h_i\) is a DSM root the receiver recognizes, \(h_i\) commits to the current fused anchor
state, and \(h_{i+1}\) is verified as the valid successor root for the transfer.

The anchor counter orders hardware events. The fused anchor head binds the hardware lineages.
DSM validity decides whether the state transition is real.

\section{Cryptographic Preliminaries}

Let \(\Hf\) denote BLAKE3 256, modeled as collision resistant and resistant to second preimages.
Let \(\HKDF\) denote a domain separated key derivation function.

Let \((\StepKeyGen,\StepSign,\StepVerify)\) be the witness signature scheme fixed by the appliance
profile. The concrete profile in this document is \WOTS\ over BLAKE3.

Let \((\PartSign,\PartVerify)\) be the RP2350 secure partition signature scheme under a partition
key generated at appliance birth and bound into the anchor bundle.

All structured objects use canonical byte encoding. If \(X\) is structured, \(\mathrm{enc}(X)\)
means its canonical byte encoding. Verifiers reject non canonical encodings.

\begin{definition}[DSM Root]
A DSM root is the sparse Merkle tree root that commits to the current local DSM state, including
relationship tips, object leaves, authority policy, anchor bundle, fused anchor head, boot head,
and anchor counter used for offline bearer mode.
\end{definition}

\begin{definition}[Spendable Parent]
A spendable parent is a DSM root \(h_i\) that commits to a spendable object, owner, relationship
context, authority policy, anchor bundle \(\Bundle\), fused anchor head \(\Ahead_i\), boot head
\(\Bhead_b\), and anchor counter \(u_i\). A valid offline bearer transfer must consume that
spendable parent into one successor root.
\end{definition}

\begin{definition}[Offline Bearer Transfer]
An offline bearer transfer is a DSM transition accepted without querying a network, storage node,
ledger, validator, sequencer, or clock service at the moment of exchange. Acceptance is decided
from canonical DSM bytes, SMT proofs, receiver challenge binding, boot ticket verification,
partition certificate verification, hardware witness verification, authenticated TROPIC01 counter
evidence, fused anchor head verification, and the receiver accepted previous root.
\end{definition}

\begin{definition}[Anchor Counter]
The anchor counter is the derived increasing DSM value:

\[
u = H_0-H.
\]

Here \(H_0\) is the TROPIC01 counter value at enrollment and \(H\) is the live TROPIC01 counter
value. TROPIC01 counters count down, so each successful counter update maps \(H\leftarrow H-1\)
and \(u\leftarrow u+1\).
\end{definition}

\section{Threat Model}

\begin{definition}[Software Clone]
A software clone receives all host readable wallet state:

\[
\mathsf{Clone}
=
(
\mathsf{seed},
\mathsf{keys},
\mathsf{chain\ history},
\mathsf{local\ database},
\mathsf{cached\ proofs},
\mathsf{host\ files},
\mathsf{application\ state}
).
\]

The clone may run on another phone, emulator, rooted host, or modified controller. It does not
receive RP2350 secure partition non exportable state or TROPIC01 internal MACANDD and counter
state.
\end{definition}

\begin{definition}[New Hardware Clone]
A new hardware clone is a device that has copied host readable state and enrollment data, but has
different RP2350 partition state, different TROPIC01 state, different partition key, different
TROPIC01 anchor identifier, different MACANDD slot state, or different physical counter state.
\end{definition}

\begin{definition}[RP2350 Partition Breach]
An RP2350 partition breach means the attacker can make arbitrary policy side calls, modify local
appliance state, or drive the software around TROPIC01 incorrectly. The protocol does not rely on
RP2350 claims alone for receiver acceptance. Such a breach may cause denial of service, counter
wasting, invalid certificates, or a bricked local anchor. It must not make an honest receiver accept
an invalid DSM transition or a second valid successor from the same spendable parent unless the
other required evidence also verifies.
\end{definition}

\begin{definition}[TROPIC01 Physical Break]
A TROPIC01 physical break means the attacker extracts or forges the secure element state needed
to produce MACANDD outputs or counterfeit counter evidence. This is outside ordinary TROPIC-only
security. In this fused design, TROPIC01 break alone is still insufficient because receiver
acceptance also requires the RP2350 partition certificate, DSM proof, boot ticket, and fused anchor
head update.
\end{definition}

\begin{definition}[Perfect Live State Clone]
A perfect live state clone is an adversary that extracts the exact current non exportable state of
the RP2350 partition, the exact current non exportable TROPIC01 MACANDD and counter state, all
DSM authority state, and can emulate those states perfectly without divergence. No offline-only
protocol can distinguish such an exact emulation from the original device.
\end{definition}

\begin{definition}[Double Spend]
A double spend exists if two distinct transitions:

\[
\tau_A:h_i\rightarrow h_A,\qquad
\tau_B:h_i\rightarrow h_B,\qquad
h_A\neq h_B,
\]

consume the same spendable parent \(h_i\) and both satisfy the DSM offline bearer acceptance
predicate for honest receivers.
\end{definition}

\begin{definition}[Closed Branch]
A closed branch is a branch created among devices controlled by the same adversary. It may be
internally consistent only inside that adversary controlled set. Since new independent
relationships require online reconciliation, the branch breaks when it meets real reconciliation
or an honest counterparty that checks the fused anchor lineage, counter evidence, and DSM root
lineage.
\end{definition}

\section{One Way Birth Fuse}

The first hardening step is to make enrollment non recreatable from public enrollment data.

\begin{definition}[One Way Birth Fuse]
The one way birth fuse \(s_{\mathrm{birth}}\) is a secret enrollment preimage formed from RP2350
partition entropy, TROPIC01 birth witness material, host entropy, device context, and authority
policy. The public anchor bundle and initial fused heads commit only to \(\Hf(s_{\mathrm{birth}})\).
The preimage \(s_{\mathrm{birth}}\) is destroyed immediately after deriving the initial private
ratchet state.
\end{definition}

At birth, the appliance samples or derives:

\[
s_{\mathrm{birth}}
=
\Hf(
  \ctx{DSM/anchor/birth-secret/v1}
  \concat \mathsf{partition\_trng}
  \concat \mathsf{tropic\_birth\_witness}
  \concat \mathsf{host\_nonce}
  \concat \mathsf{device\_id}
  \concat \mathsf{policy\_hash}
).
\]

The public birth commitment is:

\[
S_{\mathrm{birth}}=\Hf(s_{\mathrm{birth}}).
\]

The raw \(s_{\mathrm{birth}}\) is never exported.

\begin{remark}
The enrolled TROPIC01 counter value \(H_0\) is not destroyed. Receivers need \(H_0\), or a policy
pinned equivalent, to verify \(u=H_0-H\). The destroyed value is the birth preimage
\(s_{\mathrm{birth}}\), not \(H_0\).
\end{remark}

\section{Anchor Bundle}

\begin{definition}[Anchor Bundle]
The anchor bundle \(\Bundle\) is the immutable enrollment digest binding the partition public key,
partition device identifier, TROPIC01 anchor identifier, enrolled TROPIC01 counter value \(H_0\),
MACANDD boot slot, MACANDD transfer slot, DSM device identifier, authority policy, and the hash of
the destroyed one way birth fuse.
\end{definition}

Let \(q_{\mathrm{boot}}\) be the MACANDD slot used for boot fencing. Let \(q_{\mathrm{tx}}\) be the
MACANDD slot used for transfer witnesses. Then:

\[
\Bundle
=
\Hf(
  \ctx{DSM/anchor-bundle/v1}
  \concat \mathsf{partition\_pk}
  \concat \mathsf{partition\_device\_id}
  \concat \mathsf{tropic\_anchor\_id}
  \concat H_0
  \concat q_{\mathrm{boot}}
  \concat q_{\mathrm{tx}}
  \concat \mathsf{device\_id}
  \concat \mathsf{policy\_hash}
  \concat S_{\mathrm{birth}}
).
\]

Offline bearer releases under a different bundle are not valid successors of roots committed to
\(\Bundle\).

\section{Initial Fused State}

The first fused anchor head is:

\[
\Ahead_0
=
\Hf(
  \ctx{DSM/fused-anchor-head/init/v1}
  \concat \Bundle
  \concat h_0
  \concat 0
  \concat S_{\mathrm{birth}}
).
\]

The first boot head is:

\[
\Bhead_0
=
\Hf(
  \ctx{DSM/fused-boot-head/init/v1}
  \concat \Bundle
  \concat \Ahead_0
  \concat 0
  \concat S_{\mathrm{birth}}
).
\]

The initial partition ratchet is:

\[
p_0
=
\HKDF\!\left(
  \mathsf{secret}=s_{\mathrm{birth}},
  \mathsf{context}=
  \ctx{DSM/partition-ratchet-seed/v1}
  \concat \Bundle
  \concat \Ahead_0
  \concat \Bhead_0
\right).
\]

After deriving \(p_0\), the appliance destroys:

\[
s_{\mathrm{birth}}\leftarrow\bot.
\]

The device carries only forward moving private state:

\[
(
p_i,
\mathsf{TROPIC01\ MACANDD\ boot\ slot},
\mathsf{TROPIC01\ MACANDD\ transfer\ slot},
H
).
\]

The public DSM state carries:

\[
(
\Bundle,
\Ahead_i,
\Bhead_b,
u_i
).
\]

\section{Fused Anchor Head}

\begin{definition}[Fused Anchor Head]
The fused anchor head \(\Ahead_i\) is the DSM committed digest of the current offline bearer
anchor lineage. It binds the DSM root lineage, the RP2350 secure partition lineage, and the
TROPIC01 MACANDD and counter lineage into one non interchangeable state.
\end{definition}

A candidate offline successor must advance:

\[
\Ahead_i\rightarrow\Ahead_{i+1}.
\]

The next DSM root must commit to \(\Ahead_{i+1}\). A release that does not produce the next fused
anchor head is not an offline bearer successor.

\section{Boot Fenced Fused Anchor}

\begin{definition}[Boot Fenced Fused Anchor]
Offline bearer mode is enabled only after the appliance produces a boot ticket chained from the
DSM committed boot head. The boot ticket is produced internally by the firmware target from a
device authoritative boot measurement, the RP2350 secure partition boot ratchet, and a TROPIC01
boot MACANDD slot. A host request cannot drive boot head advancement and cannot supply an attacker
chosen firmware measurement. Every offline release binds the current boot ticket or boot chain. A
copied state image on new hardware cannot resume offline bearer mode because the new hardware
cannot advance the committed boot head under the enrolled anchor bundle.
\end{definition}

On boot, the firmware target advances:

\[
\Bhead_b\rightarrow\Bhead_{b+1}.
\]

The boot input is formed from device internal state:

\[
X^{\mathrm{boot}}_{b+1}
=
\Hf(
  \ctx{DSM/boot-fuse-input/v1}
  \concat \Bundle
  \concat \Ahead_i
  \concat \Bhead_b
  \concat \mathsf{boot\_seq}
  \concat \mathsf{firmware\_measurement}
  \concat \mathsf{partition\_device\_id}
).
\]

The values \(\mathsf{boot\_seq}\) and \(\mathsf{firmware\_measurement}\) are device supplied.
They are not fields accepted from the host transport request.

TROPIC01 consumes the boot MACANDD slot:

\[
W^{T,\mathrm{boot}}_{b+1}
=
\MACD(q_{\mathrm{boot}},X^{\mathrm{boot}}_{b+1}).
\]

The partition advances its boot ratchet:

\[
p_{b+1}
=
\Hf(
  \ctx{DSM/partition-boot-ratchet/v1}
  \concat p_b
  \concat W^{T,\mathrm{boot}}_{b+1}
  \concat \Bundle
  \concat \Ahead_i
  \concat \Bhead_b
  \concat \mathsf{boot\_seq}
  \concat \mathsf{firmware\_measurement}
).
\]

The old partition ratchet is erased:

\[
p_b\leftarrow\bot.
\]

The partition boot certificate message is:

\[
M^P_{\mathrm{boot},b+1}
=
\Hf(
    \ctx{DSM/partition-boot-cert/v1}
    \concat \Bundle
    \concat \Ahead_i
    \concat \Bhead_b
    \concat X^{\mathrm{boot}}_{b+1}
    \concat \Hf(W^{T,\mathrm{boot}}_{b+1})
    \concat \mathsf{boot\_seq}
    \concat \mathsf{firmware\_measurement}
).
\]

The partition signs:

\[
\sigma^P_{\mathrm{boot},b+1}=\PartSign(M^P_{\mathrm{boot},b+1}).
\]

The fixed width boot signature commitment is:

\[
\Sigma^P_{\mathrm{boot},b+1}=\SigCommit(\sigma^P_{\mathrm{boot},b+1}).
\]

The new boot head is:

\[
\Bhead_{b+1}
=
\Hf(
  \ctx{DSM/fused-boot-head/v1}
  \concat \Bundle
  \concat \Ahead_i
  \concat \Bhead_b
  \concat X^{\mathrm{boot}}_{b+1}
  \concat \Hf(p_{b+1})
  \concat \Hf(W^{T,\mathrm{boot}}_{b+1})
  \concat \Sigma^P_{\mathrm{boot},b+1}
  \concat \mathsf{firmware\_measurement}
).
\]

The boot ticket is:

\[
\BootTicket_{b+1}
=
(
\Bundle,
\Ahead_i,
\Bhead_b,
\Bhead_{b+1},
\mathsf{boot\_seq},
\mathsf{firmware\_measurement},
\sigma^P_{\mathrm{boot},b+1},
X^{\mathrm{boot}}_{b+1},
\mathsf{tropic\_boot\_witness}
).
\]

If multiple boots occur between DSM transfers, the release carries a boot chain:

\[
\BootChain =
(\BootTicket_{b+1},\ldots,\BootTicket_{b+k}),
\]

which proves:

\[
\Bhead_b\rightarrow\Bhead_{b+k}.
\]

The next DSM root commits to the latest boot head used by the release.

\section{State Bound to the Previous Root}

The key simplification is to put the fused anchor state into the DSM state committed by the
previous root.

A spendable parent root \(h_i\) commits to:

\[
(
\mathsf{object\_id},
\mathsf{owner},
\mathsf{balance\ or\ object\ state},
\mathsf{relationship\ context},
\mathsf{authority\ policy\ hash},
\Bundle,
\Ahead_i,
\Bhead_b,
u_i
).
\]

Therefore the only valid offline successor anchor counter is:

\[
u_{i+1}=u_i+1.
\]

A receiver that sees a candidate transfer from \(h_i\) rejects it unless the candidate release
proves:

\[
(h_i,\Ahead_i,\Bhead_b,u_i)
\rightarrow
(h_{i+1},\Ahead_{i+1},\Bhead_{b'},u_i+1).
\]

This removes the need for a table of offered or pending precommits. The previous root itself
carries the expected fused anchor state.

\section{TROPIC01 Counter Evidence}

The receiver must not trust a host string that says ``the counter moved.'' The receiver must also
not treat a sender supplied \(\mathsf{live\_counter}\) field as proof. That value is only a claim
unless it is authenticated by TROPIC01.

\begin{definition}[Counter Evidence]
Counter evidence for a release is evidence that the pinned TROPIC01 counter has reached the live
physical value corresponding to \(u_i+1\). It is accepted only if obtained by one of:

\begin{enumerate}[label=(\arabic*),leftmargin=2em]
\item a receiver operated authenticated \Lthree session to TROPIC01 through a verifier pairing
slot, where the host only relays encrypted packets;
\item a TROPIC01 primitive that directly authenticates the live counter value under a key pinned
by the receiver;
\item a provisioning transcript and live audit rule that is explicitly accepted by the authority
policy.
\end{enumerate}
\end{definition}

The accepted counter value is the value obtained by the receiver through an allowed counter
evidence path. The preferred path is an authenticated \Lthree verifier session. The receiver is
the endpoint and the phone or Pico only relays encrypted packets.

The receiver reads the live counter value \(H_{\mathrm{attested}}\) from TROPIC01 and checks:

\[
H_{\mathrm{attested}} = H_0-(u_i+1).
\]

Equivalently, the receiver derives:

\[
u_{\mathrm{attested}} = H_0 - H_{\mathrm{attested}}
\]

and requires:

\[
u_{\mathrm{attested}}=u_i+1.
\]

The RP2350 may relay packets, but it is not the trusted endpoint. The release may carry a counter
field for transport convenience, but that field is not trusted. The receiver accepts only the
transcript attested value, or another policy approved TROPIC01 authenticated counter value.

\begin{remark}
The preferred path is implemented and validated on target hardware. The receiver device runs its
own libtropic session end to end: the Noise handshake and the authenticated counter read are
clocked to the chip as raw SPI frames through the appliance's transparent passthrough bridge
(\texttt{OP\_SPI\_PASSTHROUGH}, Section~\ref{sec:wire}), carried device to device over the DSM
BLE transport. The sender's phone and the sender's firmware relay opaque encrypted bytes only;
neither can read or forge the attested value.
\end{remark}

\begin{remark}
If the receiver cannot obtain authenticated counter evidence, the receiver does not accept offline
bearer mode. The transfer routes to online checked mode.
\end{remark}

\section{TROPIC01 MACANDD Witness}

MACANDD is used as a hardware witness. It is not used as standalone transaction authority.

\begin{definition}[MACANDD Command Shape]
A MACANDD call is modeled as:

\[
W=\MACD(q,X),
\]

where \(q\) is a MACANDD slot index, \(X\in\bits^{256}\) is a 32 byte input, and
\(W\in\bits^{256}\) is the 32 byte output returned by TROPIC01 over an authenticated \Lthree
session.
\end{definition}

\begin{definition}[MACANDD Slot Evolution]
Let \(V_t\) be the old slot state before a MACANDD call. Let \(X\) be the input. The call computes
a new state:

\[
V_{t+1}=F_1(X\concat q),
\]

stores \(V_{t+1}\), and returns:

\[
W_t=F_2(V_t\concat V_{t+1}\concat q).
\]

The functions \(F_1\) and \(F_2\) are keyed inside TROPIC01. The host does not know their keys.
\end{definition}

\section{DSM Transition Digest}

Let \(\Delta_{i+1}\) be the canonical DSM transition package. It contains the action, recipient,
object identifier, payload, old leaf proof, new leaf proof, and all data required for the receiver
to verify:

\[
h_i\rightarrow h_{i+1}.
\]

The transition digest is:

\[
D_{i+1}=\Hf(
  \ctx{DSM/root-advance/transition-digest/v1}
  \concat \mathrm{enc}(\Delta_{i+1})
).
\]

The receiver supplies a fresh random challenge:

\[
r_R \xleftarrow{\$} \bits^{256}.
\]

The challenge is bound into the boot fenced root advance message and release. A release for one
receiver challenge is not accepted for another challenge.

\section{Boot Bound Root Advance Message}

Let \(\Bhead_{b'}\) be the current boot head proven by a boot ticket or boot chain from the DSM
committed boot head \(\Bhead_b\).

The boot bound root advance message is:

\[
M_{i+1}
=
\Hf(
  \ctx{DSM/fused-root-advance-message/v1}
  \concat \Bundle
  \concat \Ahead_i
  \concat \Bhead_{b'}
  \concat h_i
  \concat h_{i+1}
  \concat u_i
  \concat (u_i+1)
  \concat D_{i+1}
  \concat \mathsf{recipient\_device\_id}
  \concat \mathsf{object\_id}
  \concat \mathsf{authority\_policy\_hash}
  \concat r_R
).
\]

Every partition certificate and TROPIC01 transfer witness must bind this same message.

\section{Partition Commitment and TROPIC Cross Binding}

The partition first commits to the root advance message. The code authoritative partition
commitment has no partition epoch and no partition nonce. The wire certificate carries only the
fields below, and the verifier recomputes the commitment from those fields:

\[
C^P_{i+1}
=
\Hf(
  \ctx{DSM/partition-commit/v1}
  \concat \Bundle
  \concat \Ahead_i
  \concat \Bhead_{b'}
  \concat M_{i+1}
).
\]

The TROPIC01 transfer witness input includes that partition commitment:

\[
X^T_{i+1}
=
\Hf(
  \ctx{DSM/tropic-fused-transfer-input/v1}
  \concat \Bundle
  \concat \Ahead_i
  \concat \Bhead_{b'}
  \concat M_{i+1}
  \concat C^P_{i+1}
  \concat q_{\mathrm{tx}}
).
\]

TROPIC01 returns:

\[
W^T_{i+1}=\MACD(q_{\mathrm{tx}},X^T_{i+1}).
\]

The witness signing seed is:

\[
K^T_{i+1}
=
\HKDF\!\left(
  \mathsf{secret}=W^T_{i+1},
  \mathsf{context}
  =
  \ctx{DSM/tropic/fused-transfer-witness-key/v1}
  \concat X^T_{i+1}
  \concat M_{i+1}
  \concat \Bundle
  \concat \Ahead_i
\right).
\]

The witness key pair is:

\[
(\skhw,\pkhw)=\StepKeyGen(K^T_{i+1}).
\]

The public key handle is:

\[
\pkh=\Hf(\ctx{DSM/tropic/pk-hash/v1}\concat\pkhw).
\]

The TROPIC witness message is:

\[
M^T_{i+1}
=
\Hf(
  \ctx{DSM/tropic/fused-transfer-witness-message/v1}
  \concat M_{i+1}
  \concat C^P_{i+1}
  \concat X^T_{i+1}
  \concat \pkh
).
\]

The TROPIC witness signature is:

\[
\sigma^T_{i+1}=\StepSign(\skhw,M^T_{i+1}).
\]

The fixed width TROPIC signature commitment is:

\[
\Sigma^T_{i+1}=\SigCommit(\sigma^T_{i+1}).
\]

The partition final certificate binds the TROPIC witness commitment back into the partition
lineage:

\[
M^P_{i+1}
=
\Hf(
  \ctx{DSM/partition-final-cert/v1}
  \concat \Bundle
  \concat \Ahead_i
  \concat \Bhead_{b'}
  \concat M_{i+1}
  \concat C^P_{i+1}
  \concat \pkh
  \concat \Sigma^T_{i+1}
  \concat (u_i+1)
).
\]

The partition certificate is:

\[
\sigma^P_{i+1}=\PartSign(M^P_{i+1}).
\]

The fixed width partition signature commitment is:

\[
\Sigma^P_{i+1}=\SigCommit(\sigma^P_{i+1}).
\]

The cross binding is:

\[
C^P_{i+1}
\rightarrow X^T_{i+1}
\rightarrow \Sigma^T_{i+1}
\rightarrow M^P_{i+1}
\rightarrow \Sigma^P_{i+1}.
\]

Thus the TROPIC witness is bound to the partition commitment, and the partition final certificate
is bound back to the TROPIC witness commitment. Neither proof can be swapped out for a proof from
another device, another boot, another root transition, or another anchor bundle.

\section{Next Fused Anchor Head}

After the receiver obtains authenticated counter evidence, the next fused anchor head is:

\[
\Ahead_{i+1}
=
\Hf(
  \ctx{DSM/fused-anchor-head/v1}
  \concat \Bundle
  \concat \Ahead_i
  \concat \Bhead_{b'}
  \concat M_{i+1}
  \concat C^P_{i+1}
  \concat \Sigma^P_{i+1}
  \concat \pkh
  \concat \Sigma^T_{i+1}
  \concat H_{\mathrm{attested}}
).
\]

The next DSM root \(h_{i+1}\) must commit to:

\[
(
\Bundle,
\Ahead_{i+1},
\Bhead_{b'},
u_i+1
).
\]

The receiver verifies that \(h_i\) committed to the old fused anchor state and that \(h_{i+1}\)
commits to the new fused anchor state.

\section{Witness Signature Scheme}

The appliance profile uses \WOTS\ over BLAKE3. The witness key signs exactly one digest, so a
one time signature is sufficient.

\begin{definition}[\WOTS\ Parameters]
Let \(n=32\), \(w=16\), \(\ell_1=64\), \(\ell_2=3\), and \(\ell=67\). A signature is:

\[
\ell\cdot n=2144
\]

bytes. The compressed public key is 32 bytes.
\end{definition}

\begin{definition}[Chain Function]
The chain function is:

\[
F(x)=\Hf(\ctx{DSM/anchor/wots-chain/v1}\concat x).
\]

For seed \(K\) and chain index \(j\):

\[
s_j=
\HKDF(
  \mathsf{secret}=K,\;
  \mathsf{context}=\ctx{DSM/anchor/wots-sk/v1}\concat\mathrm{enc}_{16}(j)
).
\]
\end{definition}

\begin{definition}[\StepKeyGen, \StepSign, \StepVerify]
The key generation, signing, and verification algorithms are the standard Winternitz chain
construction over BLAKE3:

\[
\StepKeyGen(K)\rightarrow(\skhw,\pkhw),
\]

\[
\StepSign(\skhw,d)\rightarrow\sigma,
\]

\[
\StepVerify(\pkhw,d,\sigma)\rightarrow\{0,1\}.
\]

The secret key is the seed \(K\). It is retained only long enough to sign the release, then erased.
\end{definition}

\begin{assumption}[Witness Signature Security]
Given \(\pkhw\) and one signature on digest \(d\), no efficient adversary can produce a valid
signature on \(d'\neq d\) except with negligible probability, assuming the preimage and second
preimage resistance of \(\Hf\).
\end{assumption}

\section{Root Advance Certificate}

The root advance certificate is:

\[
\Cert_{i+1}
=
(
\Bundle,
\Ahead_i,
\Ahead_{i+1},
\Bhead_b,
\Bhead_{b'},
h_i,
h_{i+1},
u_i,
u_i+1,
D_{i+1},
M_{i+1},
C^P_{i+1},
X^T_{i+1},
\pkh,
\pkhw,
\sigma^T_{i+1},
\sigma^P_{i+1},
\mathsf{anchor\_id},
q_{\mathrm{tx}},
r_R
).
\]

The release package is:

\[
\Pkg_{i+1}
=
(
\Delta_{i+1},
\BootChain,
\Cert_{i+1},
\mathsf{counter\_evidence}_{i+1}
).
\]

\section{Compact Appliance Protocol}

The appliance has three transfer states:

\[
\Ready,\qquad \Prepared,\qquad \Committed.
\]

Boot state is separate. Offline transfer is disabled until a valid boot ticket or boot chain exists
for the current boot.

\subsection{Boot}

Boot is device internal. The host transport does not submit a boot operation, boot sequence, or
firmware measurement. On boot, the appliance:

\begin{enumerate}[label=(\arabic*),leftmargin=2em]
\item reads the active \(\Bundle,\Ahead_i,\Bhead_b,u_i\);
\item obtains the firmware and policy measurement from the firmware target;
\item computes \(X^{\mathrm{boot}}_{b+1}\);
\item consumes the TROPIC01 boot MACANDD slot;
\item advances the partition boot ratchet;
\item signs the boot certificate;
\item records \(\BootTicket_{b+1}\);
\item enables offline bearer mode only if the boot ticket verifies.
\end{enumerate}

If the boot ticket cannot be produced or verified, offline bearer mode is refused and the device
routes to online recovery.

\subsection{Prepare}

The sender proposes \(\Delta_{i+1}\), \(h_i\), and \(h_{i+1}\). The receiver checks the proposed
transfer at the human and DSM level, then supplies \(r_R\).

The appliance checks:

\[
\Active.\mathsf{status}=\Ready,
\]

\[
\Active.h=h_i,
\]

\[
\Active.\Bundle=\Bundle,
\]

\[
\Active.\Ahead=\Ahead_i,
\]

\[
\Active.u=u_i,
\]

\[
u_i=H_0-H.
\]

The expression \(H_0-H\) is computed with checked subtraction. If \(H>H_0\), the state is rejected
as a counter mismatch because a down counter cannot report a value above its enrolled value.

The appliance verifies that a boot ticket or boot chain exists from the DSM committed boot head
to the current boot head \(\Bhead_{b'}\).

It constructs \(D_{i+1}\), \(M_{i+1}\), \(C^P_{i+1}\), and \(X^T_{i+1}\). It calls MACANDD on the
transfer slot, derives the witness key, forms the TROPIC witness, forms the partition final
certificate, computes \(\Ahead_{i+1}\), and constructs \(\Cert_{i+1}\) without exporting it yet.

It writes a durable prepared record:

\[
\Active
\leftarrow
(
h_i,
\Bundle,
\Ahead_i,
\Bhead_{b'},
u_i,
\Prepared,
\Cert_{i+1},
\Delta_{i+1},
\skhw
).
\]

No counter has moved. No release has been exported.

\subsection{Commit}

The appliance forms the release package without counter evidence and stores the committed candidate
durably with:

\[
\mathsf{counter\_committed}=\mathsf{false}.
\]

Before moving the counter, the appliance re pins the live anchor counter:

\[
H_0-H=u_i.
\]

If this check fails, the operation downgrades to online recovery and does not move the counter.

If \(H=0\), the operation returns:

\[
\mathsf{EXHAUSTED\_ONLINE\_ONLY}.
\]

No release is exported or committed from an exhausted counter.

If \(H>0\), the appliance issues:

\[
\UpdateCmd.
\]

A successful counter update maps:

\[
H\leftarrow H-1.
\]

The appliance marks the committed candidate as:

\[
\mathsf{counter\_committed}=\mathsf{true}
\]

and erases \(\skhw\). The post commit physical counter value is the actual decremented reading
\(H-1\), not a value reconstructed from \(H_0-(u_i+1)\).

\subsection{Counter Evidence}

The receiver obtains counter evidence from TROPIC01. In the preferred mode, the receiver opens an
authenticated \Lthree verifier session through a verifier pairing slot, with the host relaying
encrypted packets only.

The receiver obtains \(H_{\mathrm{attested}}\) and checks:

\[
H_{\mathrm{attested}}=H_0-(u_i+1).
\]

If this check fails, the receiver rejects offline bearer acceptance.

\subsection{Emit}

The appliance may export the release only after the counter commit. The receiver attaches or
records authenticated counter evidence and verifies \(\Pkg_{i+1}\).

\subsection{Finalize}

After emitting the release, the appliance may finalize only if the active anchor counter equals
the live counter derived anchor counter:

\[
\Active.u = H_0-H.
\]

If this check fails, finalization is refused and the appliance enters online recovery.

If the check holds, finalization writes:

\[
\Active
\leftarrow
(
h_{i+1},
\Bundle,
\Ahead_{i+1},
\Bhead_{b'},
u_i+1,
\Ready,
\emptyset
).
\]

If power fails before finalize, recovery re emits the same committed release and finalizes the same
successor.

\section{Receiver Acceptance Predicate}

A receiver accepts an offline bearer transfer only if all checks hold.

\begin{definition}[Accepted Root Frontier]\label{def:frontier}
For each enrolled holder, the receiver maintains a durable adopted appliance root frontier. If a
frontier value exists for the holder, the release previous root must equal it exactly. If no
frontier value exists, the relationship is at its genesis: the receiver adopts the release's own
previous root as the expected value. Genesis adoption is sound because the authenticity of that
root is carried entirely by the other predicate checks --- the previous state commitment, the boot
chain, the partition certificate, the TROPIC witness, and above all the authenticated counter
evidence, which binds the claimed anchor counter to the live physical chip. After the receiver
accepts the release \emph{and} completes its own canonical value commit, it persists the release
next root as the new frontier for that holder. Frontier adoption is deliberately deferred until
after the canonical commit: a release that verifies but fails to commit must not move the
frontier.
\end{definition}

\begin{definition}[Boot Fenced Fused Root Advance Acceptance]
Let \(\mathsf{Accept}_{\mathrm{off}}(\Pkg_{i+1})=1\) iff:

\begin{enumerate}[label=(\arabic*),leftmargin=2em]
\item all encodings are canonical;
\item \(h_i\) is the receiver accepted previous root for the received object, under the accepted
root frontier rule of Definition~\ref{def:frontier};
\item \(h_i\) commits to \(\Bundle,\Ahead_i,\Bhead_b,u_i\);
\item the boot ticket or boot chain verifies \(\Bhead_b\rightarrow\Bhead_{b'}\);
\item the claimed next anchor counter is \(u_i+1\), using checked arithmetic;
\item the receiver challenge \(r_R\) is the challenge supplied by this receiver;
\item \(D_{i+1}\) recomputes from \(\Delta_{i+1}\);
\item \(M_{i+1}\) recomputes from the bound fields;
\item \(C^P_{i+1}\) recomputes from the partition commitment fields;
\item \(X^T_{i+1}\) recomputes from \(\Bundle,\Ahead_i,\Bhead_{b'},M_{i+1},C^P_{i+1},q_{\mathrm{tx}}\);
\item \(\pkh=\Hf(\ctx{DSM/tropic/pk-hash/v1}\concat\pkhw)\);
\item \(M^T_{i+1}\) recomputes and
\[
\StepVerify(\pkhw,M^T_{i+1},\sigma^T_{i+1})=1;
\]
\item \(M^P_{i+1}\) recomputes and
\[
\PartVerify(\mathsf{partition\_pk},M^P_{i+1},\sigma^P_{i+1})=1;
\]
\item the DSM transition proof verifies \(h_i\rightarrow h_{i+1}\);
\item the transfer gives the claimed object or value to the receiver;
\item the authority policy hash matches the previous state;
\item the receiver obtains an authenticated TROPIC01 counter value \(H_{\mathrm{attested}}\);
\item the authenticated counter value satisfies
\[
H_{\mathrm{attested}} = H_0-(u_i+1);
\]
\item equivalently, the receiver derives
\[
u_{\mathrm{attested}}=H_0-H_{\mathrm{attested}}
\]
and verifies
\[
u_{\mathrm{attested}}=u_i+1;
\]
\item \(\Ahead_{i+1}\) recomputes from the fused anchor head formula;
\item \(h_{i+1}\) commits to \(\Bundle,\Ahead_{i+1},\Bhead_{b'},u_i+1\);
\item no known firmware boundary event, physical compromise event, or policy event invalidates
the anchor.
\end{enumerate}
\end{definition}

The receiver trusts public DSM verification, the receiver challenge, the boot ticket, the partition
certificate, the hardware witness signature, the fused anchor head update, and authenticated
TROPIC01 counter evidence. The receiver does not trust host state, copied wallet files, a Pico
reported counter value, or any unauthenticated counter field carried inside the release.

\section{What Happens if RP2350 Is Breached}

The design assumes the RP2350 may fail as a trusted policy speaker. Therefore an honest receiver
does not accept any fact solely because the RP2350 says it.

If the RP2350 partition is breached alone, the attacker may:

\begin{itemize}
\item feed arbitrary roots;
\item waste MACANDD calls;
\item burn counter steps;
\item export invalid packages;
\item corrupt local state;
\item brick the appliance into online recovery.
\end{itemize}

These attacks do not become successful offline bearer transfers unless the receiver acceptance
predicate is also satisfied.

A partition certificate without a matching TROPIC01 transfer witness fails. A partition certificate
without authenticated TROPIC01 counter evidence fails. A partition certificate over an invalid DSM
transition fails.

\section{What Happens if TROPIC01 Is Broken}

TROPIC01 is not the sole authority.

If TROPIC01 is broken alone, the attacker may try to forge hardware witness output or counter
evidence. That is still not enough for an honest receiver because the release also requires:

\begin{itemize}
\item the DSM transition proof;
\item the previous root commitment to \(\Bundle,\Ahead_i,\Bhead_b,u_i\);
\item a valid boot ticket or boot chain;
\item a valid RP2350 partition certificate;
\item a valid fused anchor head update;
\item the next root commitment to \(\Bundle,\Ahead_{i+1},\Bhead_{b'},u_i+1\).
\end{itemize}

A TROPIC-only break becomes offline mode suspension and authority rotation unless the other fused
anchor predicates also fail.

\section{Why New Hardware Cannot Resume}

A copied state image can contain:

\[
\Bundle,\Ahead_i,\Bhead_b,u_i,
\mathsf{history},
\mathsf{cached\ proofs},
\mathsf{public\ enrollment\ data}.
\]

A new device cannot advance the boot head:

\[
\Bhead_b\rightarrow\Bhead_{b+1}
\]

unless it has the enrolled RP2350 partition boot ratchet and the enrolled TROPIC01 boot MACANDD
slot state. A new partition key, new partition state, new TROPIC01 anchor identifier, new MACANDD
slot state, or new physical counter gives a different lineage.

Therefore a release from new hardware fails one of:

\begin{itemize}
\item anchor bundle equality;
\item boot chain verification;
\item partition certificate verification;
\item TROPIC01 boot witness verification;
\item TROPIC01 transfer witness verification;
\item authenticated counter evidence;
\item fused anchor head recomputation;
\item next root commitment.
\end{itemize}

New hardware requires online authority rotation. It cannot continue offline from a root committed
to another bundle and fused anchor head.

\section{Power Loss Behavior}

Power may fail between any two operations.

\subsection{Before Boot Ticket Is Durable}

Offline bearer mode is disabled. Recovery must produce a valid boot ticket or route to online
checked recovery.

\subsection{After Boot Ticket Is Durable}

Offline bearer mode may proceed if the boot ticket verifies from the DSM committed boot head.
If the boot ticket is malformed, stale, or not chained from the committed boot head, offline mode
is refused.

\subsection{Before Prepared Is Durable}

No counter has moved and no release has been exported. Recovery returns to \Ready\ if anchor state
is consistent. Otherwise the appliance enters online checked recovery.

\subsection{After Prepared Is Durable}

The prepared record may complete if \(\skhw\) and the partition record are present. If required
private state is missing, the transfer cannot complete offline and must be cancelled or resolved
online. No second transfer from the same active root is allowed while the prepared record exists.

\subsection{After Release Is Durable but Before Counter Commit}

Recovery may commit the same release if the previous root, boot head, fused anchor head, and policy
still match. Since no release was exported before counter commit, no receiver has accepted an
uncommitted spend.

\subsection{After Counter Moved but Before Commit Flag Was Durable}

If the physical counter moved but \(\mathsf{counter\_committed}\) was not durably written, recovery
does not move the counter again. If the committed candidate target anchor counter already equals
the live counter derived anchor counter, recovery marks the candidate committed and re emits the
same release.

\subsection{After Counter Commit but Before Export}

The counter has moved and the release is durable. Recovery re emits the same release package. It
does not sign a new one.

\subsection{After Export but Before Finalize}

Recovery re emits the same release and finalization advances to the same \(h_{i+1}\), guarded by:

\[
\Active.u = H_0-H.
\]

\section{Recovery}

Recovery must preserve the exact committed release until it has been re emitted and finalized.
It must not erase the committed release merely because recovery found it.

The recovery rule is:

\[
\text{if a committed release exists, re emit that same release and finalize that same successor.}
\]

The appliance does not sign a new release during recovery.

\begin{enumerate}[label=(\arabic*),leftmargin=2em]
\item If the appliance is in \(\Ready\), recovery compares the stored anchor counter with the live
counter derived anchor counter. If the stored anchor counter is lower, the appliance adopts the
live value: the physical chip is the counter authority, and adopting a higher \(u\) only shrinks
the remaining budget --- it can never mint value or enable a double spend, because the physical
counter has already moved. This is the normal state after any transfer followed by a restart of a
build that does not yet persist \(\Active\) durably, since the counter is permanently below
\(H_0\) while such a build re births at \(u=0\). If the stored anchor counter is higher than the
live derived value, the appliance fails closed: a down counter cannot rise.

\item If the appliance is in \(\Prepared\), recovery checks that the previous root still equals the
active root, the boot head is valid, and the live anchor counter still equals the record anchor
counter. If the witness key and partition record are present, the prepared record may complete. If
required private state is missing, the record must be cancelled or resolved online.

\item If a committed candidate exists with \(\mathsf{counter\_committed}=\mathsf{false}\) and its
target anchor counter is \(u_{\mathrm{live}}+1\), recovery may complete the counter update only if
the candidate previous root still equals the active root and the boot/fused anchor state still
matches. This prevents recovery from burning a counter step for a stale or divergent previous root.

\item If a committed candidate exists with \(\mathsf{counter\_committed}=\mathsf{false}\) but the
live counter already equals the candidate target anchor counter, recovery marks the candidate
committed without moving the counter again. This covers the interruption case where the physical
counter moved but the committed flag was not durably written.

\item If a committed candidate exists with \(\mathsf{counter\_committed}=\mathsf{true}\), recovery
returns a re emit outcome. The record remains committed until the same release is emitted and
finalized.

\item Finalize is allowed only if the active anchor counter equals the live counter derived anchor
counter:
\[
\Active.u = H_0-H.
\]
This prevents frontier advance when local state and the physical counter disagree.
\end{enumerate}

\section{Recovery Algorithm}

\begin{lstlisting}
recover(H0, H, Active):
    live_anchor_counter = checked_sub(H0, H)
    if live_anchor_counter == ERROR:
        return COUNTER_MISMATCH

    if firmware_boundary_invalid():
        return DOWNGRADE_ONLINE

    if rmemory_map_invalid():
        return DOWNGRADE_ONLINE

    if boot_ticket_required() and not valid_boot_ticket_or_chain():
        return DOWNGRADE_ONLINE

    if Active.status == COMMITTED:
        rec = Active.record

        if rec.counter_committed == TRUE:
            if rec.next_anchor_counter != live_anchor_counter:
                return DOWNGRADE_ONLINE
            return REEMIT_COMMITTED(rec.next_root)

        if rec.counter_committed == FALSE:
            if rec.next_anchor_counter == live_anchor_counter:
                mark_counter_committed(rec)
                return REEMIT_COMMITTED(rec.next_root)

            if rec.next_anchor_counter == live_anchor_counter + 1:
                if rec.prev_root != Active.root:
                    return DOWNGRADE_ONLINE
                if rec.anchor_bundle != Active.anchor_bundle:
                    return DOWNGRADE_ONLINE
                if rec.prev_anchor_head != Active.anchor_head:
                    return DOWNGRADE_ONLINE
                counter_update()
                mark_counter_committed(rec)
                return REEMIT_COMMITTED(rec.next_root)

            return DOWNGRADE_ONLINE

    if Active.status == PREPARED:
        if Active.anchor_counter != live_anchor_counter:
            return DOWNGRADE_ONLINE

        if Active.record.prev_root != Active.root:
            return DOWNGRADE_ONLINE

        if Active.record.prev_anchor_head != Active.anchor_head:
            return DOWNGRADE_ONLINE

        if witness_key_present(Active.record) and partition_record_present(Active.record):
            return ACCEPT_PREPARED_CAN_COMPLETE

        return ONLINE_CANCEL_OR_RESOLVE

    if Active.status == READY:
        if Active.anchor_counter < live_anchor_counter:
            Active.anchor_counter = live_anchor_counter

        if Active.anchor_counter > live_anchor_counter:
            return FAIL_CLOSED

        if H == 0:
            return EXHAUSTED_ONLINE_ONLY

        return ACCEPT(Active.root)

    return DOWNGRADE_ONLINE
\end{lstlisting}

The \(\Ready\) branch adopts the live counter when the stored value is behind. The chip is the
counter authority: adoption reconciles the model to the physical position and only shrinks the
remaining spend budget. The fail closed branch is unchanged, because a stored value ahead of a
down counter is physically impossible for an honest device.

\section{Online Checked Mode and New Relationships}

New independent relationships are not created under this offline bearer authority. They require
online checked reconciliation. This matters for collusive closed branches.

If the same adversary controls both sides of an offline exchange, it can create a private branch
that only its own devices accept. That does not affect honest receivers. When the branch attempts
to meet real reconciliation or a new independent relationship, the DSM root, anchor bundle, fused
anchor head, boot head, anchor counter, and counter evidence no longer line up. The branch has
become its own reality and breaks away from the accepted one.

Therefore the meaningful security target is an honest receiver accepting value from a sender. A
closed adversary branch is not a successful attack on anyone else.

\section{Tripwire Composition}

Tripwire supplies fork exposure on reconciliation. Each device commits relationship tips into a
per device sparse Merkle tree. A valid receipt proves adjacency from an old root to a new root and
binds the relevant relationship tip update.

Tripwire does not make offline receivers instantly aware of each other. It exposes conflicting
accepted tips when they are compared.

\begin{assumption}[Tripwire Security]
Assume the hash function is collision resistant and DSM signatures are secure against chosen
message forgery. Then two distinct accepted successors from the same predecessor cannot survive
reconciliation without a hash collision, signature forgery, violated DSM predicate, or violated
hardware evidence condition.
\end{assumption}

\section{Security Claims}

\begin{theorem}[Birth Non Recreation]
Public enrollment data is insufficient to recreate the initial fused anchor lineage on new
hardware.
\end{theorem}

\begin{proof}
The anchor bundle, initial fused anchor head, initial boot head, and initial partition ratchet are
derived from \(s_{\mathrm{birth}}\), but only \(\Hf(s_{\mathrm{birth}})\) is public. The preimage
\(s_{\mathrm{birth}}\) is destroyed after deriving the initial private ratchet. A new device that
has only public enrollment data cannot derive \(p_0\), the TROPIC01 MACANDD slot state, or the
same fused anchor lineage except by inverting \(\Hf\), extracting the original non exportable
state, or perfectly emulating the original device state.
\end{proof}

\begin{theorem}[No New Hardware Resume]
Let a DSM root commit to anchor bundle \(\Bundle\), fused anchor head \(\Ahead_i\), boot head
\(\Bhead_b\), and anchor counter \(u_i\). A device with different partition hardware state or
different TROPIC01 boot state cannot produce an accepted offline bearer successor from that root,
except by forging the partition boot certificate, forging the TROPIC01 boot witness, breaking the
hash binding, or extracting and perfectly emulating the original non exportable live states.
\end{theorem}

\begin{proof}
An accepted release requires a boot ticket or boot chain from \(\Bhead_b\) to \(\Bhead_{b'}\).
The boot ticket is produced from the RP2350 partition boot ratchet and the TROPIC01 boot MACANDD
slot under the committed bundle \(\Bundle\). New hardware has a different partition state,
different TROPIC01 state, or a different bundle. Therefore it cannot advance the committed boot
head to the accepted current boot head unless it forges the required evidence or exactly emulates
the original non exportable live states.
\end{proof}

\begin{theorem}[Clone Exclusion]
A software clone cannot produce an accepted offline bearer root advance for an enrolled authority
unless it obtains the original partition evidence and TROPIC01 MACANDD output, or forges one of
the required signatures.
\end{theorem}

\begin{proof}
An accepted release requires a valid boot ticket, partition certificate, TROPIC witness signature,
authenticated TROPIC01 counter evidence, DSM transition proof, and fused anchor head update. A
software clone may copy host files and wallet state, but it does not have the RP2350 partition
ratchet, partition signing state, TROPIC01 MACANDD slot state, or live counter. Therefore it cannot
produce the accepted evidence set unless it obtains those non exportable states or forges the
required evidence.
\end{proof}

\begin{theorem}[Root Rebinding Exclusion]
A witness produced for one root advance cannot be accepted for another root advance except with
negligible probability.
\end{theorem}

\begin{proof}
The boot bound root advance message \(M_{i+1}\), partition commitment \(C^P_{i+1}\), TROPIC input
\(X^T_{i+1}\), partition certificate, TROPIC witness, and fused anchor head bind the anchor bundle,
previous fused anchor head, current boot head, previous root, next root, anchor counter, next
anchor counter, transition digest, recipient, object, policy, and receiver challenge. Changing any
of those values changes the verified messages. The old evidence no longer verifies except through
hash collision, signature forgery, or a break of the witness scheme.
\end{proof}

\begin{assumption}[One-Use Physical Counter Witness]
The offline-bearer exclusion treats the authenticated counter advance \(u_i\to u_i+1\) as a
\emph{one-use physical lineage witness}: a genuine advance of the physical counter that authorizes
exactly one accepted successor from the anchor state committing \(u_i\). A bare authenticated counter
\emph{reading} --- evidence that the live physical counter equals \(H_0-(u_i+1)\) --- is a scalar, not
a per-transition witness: it is satisfied by \emph{any} release claiming the step \(u_i\to u_i+1\), so a
single genuine advance is corroborated for two distinct same-step releases delivered to two isolated
receivers. The one-use property therefore does not follow from monotonicity alone; it is an assumption
on the anchor evidence (the offline anchor evidence assumption of the DSM formal kernel). A chip
signature over a transfer and a counter value is likewise not sufficient by itself. Where the physical
realization provides only the scalar reading, exclusion of a second same-step successor holds within a
single receiver's realized history --- the consumed parent set is threaded through one \(\Sigma\) ---
and, across isolated receivers, degrades to Tripwire exposure on reconciliation rather than local
rejection.
\end{assumption}

\begin{theorem}[Counter Step Uniqueness]
For a previous root that commits to \(u_i\), an honest receiver accepts only a release whose
authenticated TROPIC01 counter evidence proves the next anchor counter \(u_i+1\).
\end{theorem}

\begin{proof}
The receiver reads \(u_i\) from the previous DSM state and requires \(u_i+1\) as the next anchor
counter. Authenticated counter evidence must show the live TROPIC01 counter value:

\[
H_0-(u_i+1).
\]

Within a single receiver's realized history the consumed parent set records the step, so a second
successor from the same parent fails the linearity check. Across two isolated receivers the physical
counter neither increases nor resets, yet each reads the same genuine value \(H_0-(u_i+1)\); a scalar
reading is satisfied by any release claiming the step \(u_i\to u_i+1\). By the One-Use Physical Counter
Witness assumption the genuine advance authorizes exactly one accepted successor. Where only the scalar
reading is realized, a second same-step successor to a second isolated receiver is not rejected by the
local predicate and is exposed on reconciliation.
\end{proof}

\begin{theorem}[TROPIC01 Is Necessary but Not Sufficient]
A TROPIC01 witness and counter value alone do not authorize an offline bearer transfer.
\end{theorem}

\begin{proof}
The receiver acceptance predicate also requires public DSM transition validity, previous root
commitment to the fused anchor state, boot ticket verification, partition certificate verification,
receiver challenge binding, fused anchor head recomputation, and next root commitment to the new
fused anchor state. TROPIC01 evidence alone cannot satisfy those checks.
\end{proof}

\begin{theorem}[Partition Is Necessary but Not Sufficient]
An RP2350 partition certificate alone does not authorize an offline bearer transfer.
\end{theorem}

\begin{proof}
The receiver acceptance predicate also requires a TROPIC01 transfer witness, authenticated TROPIC01
counter evidence, DSM transition proof, receiver challenge binding, boot ticket verification, and
fused anchor head recomputation. A partition certificate alone cannot satisfy those checks.
\end{proof}

\begin{theorem}[Replay Idempotence]
Re emitting the same release package does not create a second spend.
\end{theorem}

\begin{proof}
The same package has the same previous root, next root, transition digest, boot chain, partition
certificate, TROPIC witness, receiver challenge, fused anchor head, and counter evidence. A
receiver that already accepted it recognizes the same transition. Re emission is duplicate
delivery, not a distinct successor.
\end{proof}

\begin{theorem}[Recoverable Commit]
If the counter moves for a release whose committed candidate is durable, recovery either re emits
the same release or downgrades to online recovery. It does not sign a different release for the
same counter step.
\end{theorem}

\begin{proof}
The recovery algorithm treats a committed record as a re emission obligation. If the committed flag
is true and the record next anchor counter matches the live counter derived anchor counter,
recovery returns \(\mathsf{REEMIT\_COMMITTED}\). If the physical counter moved but the committed
flag was not durably written, then the record target anchor counter equals the live counter derived
anchor counter, so recovery marks the same record committed and returns
\(\mathsf{REEMIT\_COMMITTED}\). If the counter has not moved and the target anchor counter is the
next physical step, recovery may move the counter only after checking that the committed candidate
previous root, bundle, and fused anchor head still equal the active state. In no branch does
recovery derive a new witness key or sign a different release.
\end{proof}

\begin{theorem}[No Accepted Offline Bearer Double Spend]
Under the stated assumptions, no adversary can produce two distinct accepted offline bearer
successors from the same spendable parent except with negligible probability. If a conflicting
branch is created outside the model, it is exposed on reconciliation.
\end{theorem}

\begin{proof}
For a previous root \(h_i\), the previous state commits to anchor bundle \(\Bundle\), fused anchor
head \(\Ahead_i\), boot head \(\Bhead_b\), and anchor counter \(u_i\). An honest receiver accepts
only a transition to \(h_{i+1}\) with next anchor counter \(u_i+1\), valid DSM proof, valid boot
ticket, valid partition certificate, valid TROPIC witness, authenticated TROPIC01 counter evidence,
valid fused anchor head update, and next root commitment to the new fused anchor state.

A software clone cannot produce the partition and TROPIC evidence. A new hardware clone cannot
resume the committed boot head. A broken TROPIC01 alone cannot fake the partition certificate, DSM
proof, boot ticket, and fused anchor head update. A breached RP2350 alone cannot forge counter
evidence --- but it does not need to: a single genuine advance to \(u_i+1\) is a scalar corroborated
for every same-step release, so exclusion of a second accepted package from the same anchor counter
rests on the One-Use Physical Counter Witness assumption, not on monotonicity. Under that assumption
the genuine advance authorizes exactly one accepted successor, and a different successor is rejected by
the receiver predicate. Where the realization provides only the scalar counter reading and the RP2350
partition is breached --- so the partition certificate and witness are attacker-produced --- two
distinct same-step successors delivered to two isolated receivers each satisfy the local predicate; the
second is not rejected locally and is exposed on reconciliation. The exclusion is then a property of a
single reconciled realized history against a live consumed parent set, not of the isolated per-receiver
predicate. This is the boundary of the offline claim: prevention while the one-use witness holds and the
partition is uncompromised; Tripwire detection and attribution on reconciliation otherwise.
\end{proof}

\section{TLA\texorpdfstring{\textsuperscript{+}}{+} Model}

\begin{lstlisting}[language=TLA]
VARIABLES H, Active, CurrentRoot, step, delivered

LiveAnchorCounter == H0 - H
NextAnchorCounter == LiveAnchorCounter + 1

OfflineReady ==
    /\ Active.root = CurrentRoot
    /\ Active.anchor_counter = LiveAnchorCounter
    /\ Active.status = "ready"
    /\ Active.boot_valid = TRUE
    /\ H > 0
    /\ RMemoryMapOk
    /\ FirmwareBoundaryOk
    /\ LockoutOk

Boot ==
    /\ step = "boot_required"
    /\ BootWitnessGenerated
    /\ PartitionBootCertValid
    /\ Active' = [Active EXCEPT !.boot_head = NextBootHead,
                                !.boot_valid = TRUE]
    /\ step' = "idle"
    /\ UNCHANGED <<H, CurrentRoot, delivered>>

Prepare ==
    /\ step = "idle"
    /\ OfflineReady
    /\ ReceiverChallengeFresh
    /\ TransitionDigestOk
    /\ PartitionCommitGenerated
    /\ TropicWitnessGenerated
    /\ PartitionFinalCertGenerated
    /\ NextAnchorHeadComputed
    /\ Active' = [
         root           |-> Active.root,
         anchor_bundle  |-> Active.anchor_bundle,
         anchor_head    |-> Active.anchor_head,
         boot_head      |-> Active.boot_head,
         anchor_counter |-> Active.anchor_counter,
         boot_valid     |-> Active.boot_valid,
         status         |-> "prepared",
         record         |-> [
             prev_root           |-> Active.root,
             next_root           |-> NextRoot,
             prev_anchor_head    |-> Active.anchor_head,
             next_anchor_head    |-> NextAnchorHead,
             boot_head           |-> Active.boot_head,
             anchor_counter      |-> LiveAnchorCounter,
             next_anchor_counter |-> NextAnchorCounter,
             challenge           |-> ReceiverChallenge,
             release             |-> ReleasePkg,
             committed           |-> FALSE
         ]
       ]
    /\ step' = "prepared"
    /\ UNCHANGED <<H, CurrentRoot, delivered>>

CommitStart ==
    /\ step = "prepared"
    /\ Active.status = "prepared"
    /\ Active.record.prev_root = Active.root
    /\ Active.record.prev_anchor_head = Active.anchor_head
    /\ Active.anchor_counter = LiveAnchorCounter
    /\ ReleaseDurable
    /\ Active' = [
         root           |-> Active.root,
         anchor_bundle  |-> Active.anchor_bundle,
         anchor_head    |-> Active.anchor_head,
         boot_head      |-> Active.boot_head,
         anchor_counter |-> Active.anchor_counter,
         boot_valid     |-> Active.boot_valid,
         status         |-> "committed",
         record         |-> Active.record
       ]
    /\ step' = "commit_started"
    /\ UNCHANGED <<H, CurrentRoot, delivered>>

CommitCounter ==
    /\ step = "commit_started"
    /\ Active.status = "committed"
    /\ Active.record.committed = FALSE
    /\ Active.record.next_anchor_counter = LiveAnchorCounter + 1
    /\ Active.record.prev_root = Active.root
    /\ Active.record.prev_anchor_head = Active.anchor_head
    /\ H > 0
    /\ H' = H - 1
    /\ Active' = [Active EXCEPT !.record.committed = TRUE,
                                !.anchor_counter =
                                  Active.record.next_anchor_counter]
    /\ step' = "committed"
    /\ UNCHANGED <<CurrentRoot, delivered>>

Emit ==
    /\ step = "committed"
    /\ Active.status = "committed"
    /\ Active.record.committed = TRUE
    /\ Active.anchor_counter = LiveAnchorCounter
    /\ delivered' = delivered \cup {Active.record.release}
    /\ step' = "emitted"
    /\ UNCHANGED <<H, Active, CurrentRoot>>

Finalize ==
    /\ step \in {"committed", "emitted"}
    /\ Active.status = "committed"
    /\ Active.record.committed = TRUE
    /\ Active.anchor_counter = LiveAnchorCounter
    /\ Active' = [
         root           |-> Active.record.next_root,
         anchor_bundle  |-> Active.anchor_bundle,
         anchor_head    |-> Active.record.next_anchor_head,
         boot_head      |-> Active.record.boot_head,
         anchor_counter |-> Active.record.next_anchor_counter,
         boot_valid     |-> Active.boot_valid,
         status         |-> "ready",
         record         |-> EmptyRec
       ]
    /\ CurrentRoot' = Active.record.next_root
    /\ step' = "idle"
    /\ UNCHANGED <<H, delivered>>

PowerLoss ==
    /\ step' = "recover"
    /\ UNCHANGED <<H, Active, CurrentRoot, delivered>>

RecoverCommitted ==
    /\ step = "recover"
    /\ Active.status = "committed"
    /\ Active.record.committed = TRUE
    /\ Active.record.next_anchor_counter = LiveAnchorCounter
    /\ delivered' = delivered \cup {Active.record.release}
    /\ Active' = Active
    /\ CurrentRoot' = CurrentRoot
    /\ step' = "emitted"
    /\ UNCHANGED H

RecoverCommitFlagLost ==
    /\ step = "recover"
    /\ Active.status = "committed"
    /\ Active.record.committed = FALSE
    /\ Active.record.next_anchor_counter = LiveAnchorCounter
    /\ Active' = [Active EXCEPT !.record.committed = TRUE,
                                !.anchor_counter =
                                  Active.record.next_anchor_counter]
    /\ delivered' = delivered \cup {Active.record.release}
    /\ CurrentRoot' = CurrentRoot
    /\ step' = "emitted"
    /\ UNCHANGED H

RecoverCommitNotMoved ==
    /\ step = "recover"
    /\ Active.status = "committed"
    /\ Active.record.committed = FALSE
    /\ Active.record.next_anchor_counter = LiveAnchorCounter + 1
    /\ Active.record.prev_root = Active.root
    /\ Active.record.prev_anchor_head = Active.anchor_head
    /\ H > 0
    /\ H' = H - 1
    /\ Active' = [Active EXCEPT !.record.committed = TRUE,
                                !.anchor_counter =
                                  Active.record.next_anchor_counter]
    /\ delivered' = delivered \cup {Active.record.release}
    /\ CurrentRoot' = CurrentRoot
    /\ step' = "emitted"

RecoverPrepared ==
    /\ step = "recover"
    /\ Active.status = "prepared"
    /\ Active.anchor_counter = LiveAnchorCounter
    /\ Active.record.prev_root = Active.root
    /\ Active.record.prev_anchor_head = Active.anchor_head
    /\ Active' = Active
    /\ delivered' = delivered
    /\ CurrentRoot' = CurrentRoot
    /\ step' = "prepared"
    /\ UNCHANGED H

RecoverIdle ==
    /\ step = "recover"
    /\ Active.status = "ready"
    /\ Active.anchor_counter = LiveAnchorCounter
    /\ Active' = Active
    /\ delivered' = delivered
    /\ CurrentRoot' = CurrentRoot
    /\ step' = "idle"
    /\ UNCHANGED H

RecoverIdleAdopt ==
    /\ step = "recover"
    /\ Active.status = "ready"
    /\ Active.anchor_counter < LiveAnchorCounter
    /\ Active' = [Active EXCEPT !.anchor_counter = LiveAnchorCounter]
    /\ delivered' = delivered
    /\ CurrentRoot' = CurrentRoot
    /\ step' = "idle"
    /\ UNCHANGED H

OnlineAdvance ==
    /\ step = "idle"
    /\ CurrentRoot' = OnlineSuccessor(CurrentRoot)
    /\ step' = "idle"
    /\ UNCHANGED <<H, Active, delivered>>

Next ==
    Boot
    \/ Prepare
    \/ CommitStart
    \/ CommitCounter
    \/ Emit
    \/ Finalize
    \/ PowerLoss
    \/ RecoverCommitted
    \/ RecoverCommitFlagLost
    \/ RecoverCommitNotMoved
    \/ RecoverPrepared
    \/ RecoverIdle
    \/ RecoverIdleAdopt
    \/ OnlineAdvance
\end{lstlisting}

The model obligations are:

\begin{description}[leftmargin=1.8em,style=nextline]
\item[Single active root.]
The appliance has one active root and one active anchor counter.

\item[One fused anchor head.]
The active root commits to one fused anchor head.

\item[Boot before offline.]
Offline transfer is disabled unless a boot ticket or boot chain verifies from the DSM committed
boot head.

\item[No export before commit.]
A release is emitted only after the counter moves.

\item[Replay is idempotent.]
The same committed release may be re emitted without creating a new successor.

\item[Prepared state blocks replacement.]
A prepared record blocks another preparation from the same active root.

\item[Recovery repeats the same successor.]
Recovery emits the same committed successor. It does not sign a new one.

\item[Chip counter authority on recovery.]
A \(\Ready\) appliance whose stored anchor counter is behind the live counter derived value adopts
the live value (\(\mathsf{RecoverIdleAdopt}\)); adoption only shrinks the remaining budget. An
appliance ahead of the live value fails closed.

\item[Finalize guard.]
Finalize requires \(\Active.u=H_0-H\), equivalently
\(\Active.\mathsf{anchor\_counter}=H_0-H\).

\item[Online separation.]
If online state advances without the anchor, offline bearer mode is refused until resync.
\end{description}

\section{Wire Protocol}\label{sec:wire}

The transport uses protocol buffers. The secure core exposes one entry point:

\[
\Handle:\mathsf{bytes}\rightarrow\mathsf{bytes}.
\]

The host may relay packets, but receiver acceptance is based on public proofs, boot evidence,
partition evidence, fused anchor head verification, and authenticated TROPIC01 counter evidence.
Boot fencing is device internal. The host transport does not expose a callable boot operation, and
it does not accept host supplied boot sequence or firmware measurement fields.

\begin{lstlisting}[language=proto]
syntax = "proto3";
package dsm.anchor;

message TransitionPackage {
  bytes relationship_id = 1;
  bytes object_id = 2;
  bytes sender_device_id = 3;
  bytes recipient_device_id = 4;
  bytes prev_root = 5;
  bytes next_root = 6;
  uint64 anchor_counter = 7;
  uint64 next_anchor_counter = 8;
  uint32 action_type = 9;
  bytes action_fields = 10;
  bytes payload_hash = 11;
  bytes old_leaf_proof = 12;
  bytes new_leaf_proof = 13;
  bytes authority_policy_hash = 14;
}

message BootTicket {
  bytes anchor_bundle = 1;
  bytes anchor_head = 2;
  bytes prev_boot_head = 3;
  bytes next_boot_head = 4;
  uint64 boot_seq = 5;
  bytes firmware_measurement = 6;
  bytes partition_boot_signature = 7;
  bytes tropic_boot_input = 8;
  bytes tropic_boot_witness = 9;
}

message RootAdvanceCertificate {
  bytes anchor_bundle = 1;
  bytes prev_anchor_head = 2;
  bytes next_anchor_head = 3;
  bytes prev_boot_head = 4;
  bytes current_boot_head = 5;
  bytes prev_root = 6;
  bytes next_root = 7;
  uint64 anchor_counter = 8;
  uint64 next_anchor_counter = 9;
  bytes transition_digest = 10;
  bytes root_advance_message = 11;
  bytes partition_commitment = 12;
  bytes tropic_transfer_input = 13;
  bytes pk_hash = 14;
  bytes pk_hw = 15;
  bytes sigma_tropic = 16;
  bytes sigma_partition = 17;
  bytes anchor_id = 18;
  uint32 transfer_slot = 19;
  bytes receiver_challenge = 20;
}

message CounterEvidence {
  bytes anchor_id = 1;
  uint64 enrolled_counter = 2;
  uint64 live_counter_claim = 3;
  uint64 derived_anchor_counter_claim = 4;
  bytes verifier_transcript = 5;
}

message OfflineRelease {
  TransitionPackage transition = 1;
  repeated BootTicket boot_chain = 2;
  RootAdvanceCertificate cert = 3;
  CounterEvidence counter = 4;
}

enum Op {
  OP_UNSPECIFIED = 0;
  reserved 1;
  reserved "OP_BOOT";
  OP_PREPARE = 2;
  OP_COMMIT = 3;
  OP_EMIT = 4;
  OP_FINALIZE = 5;
  OP_STATUS = 6;
  OP_CANCEL = 7;
  OP_SPI_PASSTHROUGH = 8;
}

message ApplianceRequest {
  Op op = 1;
  TransitionPackage transition = 2;
  bytes receiver_challenge = 3;
  reserved 4, 5;
  bytes spi_payload = 6;
}

message ApplianceResponse {
  Op op = 1;
  bool ok = 2;
  uint32 error = 3;
  OfflineRelease release = 4;
  bytes active_root = 5;
  bytes anchor_bundle = 6;
  bytes active_anchor_head = 7;
  bytes active_boot_head = 8;
  uint64 active_anchor_counter = 9;
  uint32 status = 10;
  bool boot_valid = 11;
  bytes spi_response = 12;
  bytes active_committed_boot_head = 13;
  bytes pk_anchor_id = 14;
  uint64 pk_enrolled_counter = 15;
  bytes pk_partition_pk = 16;
}
\end{lstlisting}

The enum value 1 and the name \texttt{OP\_BOOT} are reserved, deliberately, so old host driven
boot semantics do not reappear. \texttt{ApplianceRequest} fields 4 and 5 are also reserved. They
must not be reused for host supplied \texttt{boot\_seq} or \texttt{firmware\_measurement}. Those
values are device internal boot data recorded in the \texttt{BootTicket}.

\texttt{OP\_SPI\_PASSTHROUGH} is not an appliance operation. It is a transparent raw SPI bridge:
\texttt{spi\_payload} carries opaque MOSI bytes that the firmware clocks to TROPIC01 unmodified,
and \texttt{spi\_response} returns the raw MISO bytes. This is the transport for the receiver
operated counter evidence path: the remote receiver runs its own authenticated libtropic session
end to end against the chip, and this device is a wire. The firmware never interprets, buffers,
or replays passthrough content.

Because TROPIC01 maintains a single authenticated \Lthree session, a receiver's passthrough
handshake re keys the chip and invalidates the appliance's own session. The firmware therefore
marks its session stale after any passthrough traffic and re establishes a fresh authenticated
session, with a fresh ephemeral key, before executing the next appliance operation. Without this
rule, the first appliance operation after a relayed counter read fails on a dead session.

The \texttt{OP\_STATUS} response fields 13 through 16 let a host client assemble the next
transition and the receiver pin material from a single status round trip: the committed boot head
(what the previous fused anchor state leaf commits), the enrolled anchor identity, the enrolled
counter value \(H_0\), and the partition public key. These are the chip self describing the
identity a receiver pins; they are convenience transport, not trusted acceptance inputs.

The partition commitment in \texttt{RootAdvanceCertificate} is recomputed from
\(\Bundle,\Ahead_i,\Bhead_{b'}\), and \(M_{i+1}\) only. It does not carry and does not
hash a partition epoch or partition nonce.

The fields \texttt{live\_counter\_claim} and \texttt{derived\_anchor\_counter\_claim} are
transport claims. They are not accepted as proof. The receiver verifier must parse or obtain an
authenticated TROPIC01 value from \texttt{verifier\_transcript}, or through another policy
approved authenticated counter evidence path.

The signature fields \texttt{sigma\_tropic} and \texttt{sigma\_partition} are variable length.
Whenever either signature is bound into another digest, the implementation uses
\(\SigCommit(\sigma)\), not raw concatenation and not a bare \(\Hf(\sigma)\).

\section{Reference Implementation}

The reference implementation has three visible parts:

\begin{itemize}
\item \texttt{crates/dsm-anchor-core}, the Rust protocol core;
\item \texttt{crates/dsm-anchor-pico}, the firmware target;
\item the DSM SDK receiver adapter that routes offline bearer acceptance through the anchor
predicate.
\end{itemize}

The protocol core implements:

\begin{itemize}
\item canonical encodings;
\item anchor bundle construction;
\item birth fuse commitment;
\item boot ticket verification;
\item root advance digest construction;
\item partition commitment construction without partition epoch or partition nonce;
\item TROPIC01 witness input construction;
\item \WOTS\ witness signatures;
\item BLAKE3-SPHINCS\textsuperscript{+} SPX128f partition signature handling;
\item fixed width \(\SigCommit(\sigma)\) binding for variable length signatures;
\item partition certificate verification;
\item fused anchor head construction;
\item public receiver acceptance;
\item recovery, with live counter adoption for a \(\Ready\) appliance behind the chip;
\item protocol buffer wire encoding.
\end{itemize}

The firmware target is the Rust \texttt{dsm-anchor-pico} crate. It drives TROPIC01 through
\texttt{libtropic-rs} bindings over SPI at 3.3 V under an authenticated \Lthree\ session. The
protocol does not require a missing C firmware layer. The RP2350 code is a transport and partition
state machine layer. Receiver acceptance is not based on trusting RP2350 statements. The firmware
implements the single session recovery rule of Section~\ref{sec:wire}: after any raw SPI
passthrough traffic it re establishes its own authenticated session before the next appliance
operation.

Boot fencing is performed by the device. The host operation value formerly used for boot is
reserved. The host does not provide the boot sequence or firmware measurement used by the boot
ratchet.

The receiver acceptance implementation separates the counter evidence parser from the release
fields. The counter verifier returns the authenticated counter value obtained from TROPIC01, and
the acceptance predicate compares that value to the expected \(H_0-(u_i+1)\).

The receiver adapter implements the accepted root frontier of Definition~\ref{def:frontier} as a
durable per holder table. The acceptance entry point takes the adopted frontier as an optional
input --- absent means relationship genesis, adopt the release's own previous root --- and returns
the adopted successor state \((\mathsf{next\_root},u_i+1)\), which the adapter persists only after
its own canonical value commit succeeds. The receiver's mirror of the sender's per step signing
lineage follows the same commit gated discipline: nothing the receiver stores about the sender
advances on a release that did not commit.

Verifier trait names should follow the same semantics:

\begin{itemize}
\item \texttt{prev\_root\_commits\_anchor\_state};
\item \texttt{verify\_transition(prev\_root, next\_root, ...)};
\item \texttt{verify\_boot\_chain};
\item \texttt{verify\_partition\_certificate};
\item \texttt{read\_authentic\_counter};
\item \texttt{next\_root\_commits\_anchor\_state}.
\end{itemize}

\section{Implementation Status}

The host side appliance state machine, fused root advance, receiver predicate, canonical wire
validation, and Pico firmware target are implemented in the reference code. Host tests are green
for the implemented anchor core behavior.

Receiver side offline bearer acceptance is live end to end on target hardware. On 5 July 2026 the
complete path was exercised between two production Android handsets, with the appliance (Pico 2 W
and TROPIC01) attached to the sender over USB and the handsets connected over BLE. All previously
listed producer inputs were present together: the sender carried a real \texttt{OfflineRelease} on
the transfer confirmation; the DSM SMT anchor state leaf committed the anchor tuple
\((\Bundle,\Ahead_i,\Bhead_b,u_i)\), reconciled to the live chip state; and the receiver operated
its own authenticated \Lthree\ verifier session through the raw SPI passthrough bridge and
obtained the attested counter. The full acceptance predicate passed every check, the anchor
counter advanced by exactly one, both sides committed canonically, and value was delivered. The
receiver adapter retains its fail closed posture: an absent or incomplete pin, an absent counter
reader, or a failed authenticated read still routes the transfer to online checked recovery.

Device side durable recovery is specified by the protocol and implemented in host logic, but it is
not yet a hardware backed cross power cycle property unless \(\Active\) is persisted in TROPIC01
R memory. A firmware build that keeps \(\Active\) in RAM and re enrolls on boot is an early
bringup build, not a complete hardware backed recovery implementation. The recovery rule for a
\(\Ready\) appliance behind the live counter --- adopt the live value --- makes that bringup mode
safe in practice: a re birth on an already provisioned chip reconciles to the true physical
counter position instead of failing on the mismatch, and adoption can only shrink the remaining
budget.

Firmware measurement is specified by the protocol, but a fixed test constant is only a bringup
stub. A production measured boot claim requires replacing that constant with a real measurement of
firmware and policy state.

The TLA\texorpdfstring{\textsuperscript{+}}{+} model in this paper is the boot fenced fused
anchor model, including the \(\mathsf{RecoverIdleAdopt}\) action matching the implemented
recovery rule. It should exist in the repository as
\texttt{tla/DSM\_BootFencedFusedAnchor.tla}; at the time of writing the model file has not yet
landed there, and the listing in this paper is the normative model. Older offline finality or
Tripwire models do not replace this model because they use different variables and prove a
different machine.

One residual is known and documented in the Limits section: the physical counter and appliance
root advance when a release is produced, even if that particular receiver ultimately rejects the
confirm, so a rejected attempt opens a gap between the sender lineage and that receiver's adopted
frontier. An online re admission flow for that receiver is specified as future work; a first
transfer to any receiver is unaffected because of genesis adoption.

\section{Validation Plan}

\begin{description}[leftmargin=1.8em,style=nextline]
\item[T1. Hardware bringup.]
Connect Raspberry Pi Pico 2 W to Secure Tropic Click over SPI at 3.3 V and confirm TROPIC01
communication.

\item[T2. TROPIC identity.]
Read chip identity and bind \(\mathsf{anchor\_id}\) to the authority policy.

\item[T3. Secure session.]
Establish authenticated \Lthree and confirm encrypted communication.

\item[T4. Birth fuse.]
Generate \(s_{\mathrm{birth}}\), derive \(\Bundle,\Ahead_0,\Bhead_0,p_0\), destroy
\(s_{\mathrm{birth}}\), and verify public enrollment data cannot recreate \(p_0\).

\item[T5. Boot ticket.]
Produce a boot ticket from the partition boot ratchet and TROPIC01 boot MACANDD slot. Expected
result: offline bearer mode is disabled until the ticket verifies.

\item[T6. New hardware resume.]
Copy host state to a different RP2350/TROPIC01 pair. Expected result: boot ticket cannot be
chained from the committed boot head under the enrolled bundle.

\item[T7. Counter direction.]
Initialize counter at \(H_0\), update it, and confirm \(u=H_0-H\) increases by one.

\item[T8. Counter evidence.]
Have the receiver act as verifier endpoint and read the live counter through an authenticated
\Lthree session.

\item[T9. MACANDD transfer witness.]
Run MACANDD on the fused transfer input and derive a witness key.

\item[T10. Public witness verification.]
Produce a release and verify the TROPIC witness signature under the public acceptance predicate.

\item[T11. Partition cross binding.]
Mutate the partition commitment, TROPIC witness, or partition final certificate. Expected result:
receiver rejects.

\item[T12. DSM proof verification.]
Mutate old leaf proof, new leaf proof, recipient, payload, object, previous root, or next root.
Expected result: receiver rejects.

\item[T13. Fused anchor head verification.]
Mutate \(\Ahead_i\), \(\Ahead_{i+1}\), \(\Bhead_b\), or \(\Bhead_{b'}\). Expected result: receiver
rejects.

\item[T14. Counter mismatch.]
Attempt to reuse the same previous root after one counter commit. Expected result: receiver rejects
because the previous state committed anchor counter and authenticated TROPIC counter evidence do
not match.

\item[T15. RP2350 breach simulation.]
Feed arbitrary next roots, wrong previous roots, stale anchor counters, forged boot heads, and
forged status outputs. Expected result: receiver rejects unless DSM proof, partition evidence,
TROPIC witness, counter evidence, and fused anchor head are valid.

\item[T16. Counter trust seam.]
Create a release with a forged host supplied counter claim but a transcript attested counter value
that does not match. Expected result: receiver rejects. The accepted value is the authenticated
TROPIC01 value, not the host claim.

\item[T17. Power loss after boot ticket.]
Cut power after boot ticket generation. Expected result: boot chain either verifies or offline
mode remains disabled.

\item[T18. Power loss after prepared.]
Cut power after prepared record is durable. Expected result: complete the same record, cancel it,
or route online. Do not create a second record from the same active root.

\item[T19. Power loss after release durable before counter commit.]
Cut power after the committed candidate is durable but before the counter moves. Expected result:
recovery may commit the same release only if the candidate previous root and fused anchor state
still equal the active state.

\item[T20. Power loss after counter moved before commit flag.]
Cut power after the physical counter moves but before the committed flag is durable. Expected
result: recovery marks the same candidate committed, does not move the counter again, and re emits
the same release.

\item[T21. Replay.]
Emit the same release twice. Expected result: duplicate delivery of the same successor.

\item[T22. Online stale.]
Advance DSM online without the anchor. Expected result: offline bearer mode refused until resync.

\item[T23. Physical compromise policy.]
Mark anchor physically suspect. Expected result: offline mode suspended and authority rotation
required.
\end{description}

\section{Validation Results on Target Hardware}

The hardware bringup and the end to end activation exercised the target stack:

\begin{center}
\begin{tabular}{lll}
\toprule
\textbf{Stage} & \textbf{Result} & \textbf{Meaning}\\
\midrule
Chip identity &
TROPIC01 identity read; silicon revision \texttt{ACAB} &
anchor can be pinned\\
Secure session &
X25519 handshake, encrypted \Lthree, echo and TRNG checked &
authenticated command path works\\
Counter &
counter initialized at \(H_0=1000\), updated to \(H=997\) &
\(u=H_0-H\) advances\\
MACANDD &
reproducible witness after rearm, distinct output without rearm &
stateful witness behavior confirmed\\
Witness flow &
\(\MACD\rightarrow K\rightarrow\WOTS\), 32 byte public key, 2144 byte signature &
public witness path works\\
Core acceptance &
release verified and successor frontier matched in host validation &
public predicate accepts the constructed release when required evidence is present\\
USB host &
external host drove status and offer flow over USB CDC &
appliance can be driven externally\\
Counter evidence over relay &
receiver handset ran its own authenticated \Lthree session through the raw SPI passthrough
bridge over BLE and read the live counter &
receiver operated counter evidence works on target hardware (T8)\\
End to end acceptance &
two handset offline bearer transfer accepted under the full predicate; anchor counter advanced by
one; canonical commit and value delivery on both sides (5 July 2026) &
the offline bearer path is live on target hardware\\
\bottomrule
\end{tabular}
\end{center}

\subsection{End to End Activation Run, 5 July 2026}

The end to end activation was a real value transfer between two production Android handsets
(Samsung SM-A165M class devices), with the appliance --- a Raspberry Pi Pico 2 W carrying the
TROPIC01 click board --- attached to the sender's handset over USB CDC, and the two handsets
connected only over BLE. No network path participated in acceptance. The concrete observed facts
of the accepted run:

\begin{itemize}
\item The sender's appliance reconciled its anchor state leaf to the live chip before the
transfer: the DSM device root committed \((\Bundle,\Ahead_i,\Bhead_b,u_i)\) with \(u_i=12\),
sourced from the chip's own status rather than any host record. Earlier rejected attempts and
appliance restarts had advanced the physical counter past the host's view; the chip as counter
authority (and, on recovery, the adopt rule of the \(\Ready\) branch) reconciled every one of
them without a counter mismatch.
\item The appliance drove \(\mathsf{PREPARE}\rightarrow\mathsf{COMMIT}\rightarrow\mathsf{EMIT}\)
against the live chip and produced a real \texttt{OfflineRelease}: MACANDD transfer witness,
\WOTS\ witness signature (2144 bytes), partition certificate, and boot chain, all bound to the
receiver's fresh challenge \(r_R\). The physical counter moved by exactly one step
(\(u:12\rightarrow 13\); on the enrolled bench chip with \(H_0=1000\), a live reading of
\(H=987\)).
\item The confirm message carrying the release together with the DSM receipts totaled roughly
221 kilobytes on the wire, dominated by post quantum signature material, and was chunked over
BLE GATT without loss under the transport retry rules.
\item The receiver ran its own authenticated \Lthree\ verifier session end to end --- X25519
handshake and authenticated counter read --- as raw SPI frames through the
\texttt{OP\_SPI\_PASSTHROUGH} bridge, relayed device to device over BLE. Individual relay round
trips were observed at roughly 150 milliseconds. The sender's handset and firmware relayed
opaque encrypted bytes only. The receiver verified
\(H_{\mathrm{attested}}=H_0-(u_i+1)=987\).
\item The full acceptance predicate passed every check. This was the first transfer from this
holder, so the accepted root frontier was adopted at genesis and persisted, with
\(\mathsf{next\_anchor\_counter}=13\), only after the receiver's canonical commit succeeded.
\item Both sides committed canonically and a 42 unit token balance moved from sender to receiver
(sender \(100\rightarrow 58\), receiver \(0\rightarrow 42\)), with the sender's own signing
lineage advanced only at commit time under the same deferred discipline.
\end{itemize}

Relative to the validation plan, this run discharges T1 through T3 and T7 through T10 live on
target hardware, and executes the boot fence gating and boot chain verification of T5 inside the
accepted release. The adversarial mutation cases (T11 through T16) remain discharged by host
tests against the same predicate implementation; the power loss cases (T17 through T20) are
specified and host validated but await hardware backed durable \(\Active\) persistence to be
meaningful across a physical power cycle.

For this boot fenced fused anchor version, the remaining work is TROPIC01 R memory backed durable
\(\Active\) persistence, replacement of the firmware measurement test constant with a real
measurement, landing the boot fenced fused anchor TLA\texorpdfstring{\textsuperscript{+}}{+}
model file in the repository, and the online re admission flow for a receiver whose adopted
frontier has fallen behind the sender lineage after a rejected attempt.

\section{Security Summary}

\begin{center}
\begin{tabular}{lll}
\toprule
\textbf{Part} & \textbf{Provides} & \textbf{Does not provide}\\
\midrule
Birth fuse & non recreatable enrollment preimage & live clone detection by itself\\
Anchor bundle & immutable hardware and policy binding & forward motion by itself\\
Boot ticket & new hardware resume resistance & DSM validity\\
Previous DSM root & object state and expected fused anchor state & hardware presence\\
Receiver challenge & freshness and recipient binding & uniqueness by itself\\
DSM transition proof & valid root update & hardware presence\\
RP2350 partition cert & appliance partition lineage & DSM validity by itself\\
TROPIC01 MACANDD & enrolled hardware witness & DSM validity\\
TROPIC01 counter evidence & exact next physical counter state & recipient binding\\
Fused anchor head & binds DSM, partition, and TROPIC lineage & value validity by itself\\
Accepted root frontier & receiver side lineage continuity & authenticity by itself at genesis\\
RP2350 host path & transport and local state machine & trusted receiver facts\\
Recovery record & no orphaned committed release & new authority by itself\\
Tripwire & fork exposure on reconciliation & instant awareness\\
 & cryptographic security\\
\bottomrule
\end{tabular}
\end{center}

\section{Limits}

\begin{enumerate}[label=(\arabic*),leftmargin=2em]
\item \textbf{Perfect live state emulation is outside offline distinguishability.}
If an attacker extracts and perfectly emulates the exact current non exportable state of the
partition, TROPIC01, and DSM authority state, no offline-only protocol can distinguish that from
the original.

\item \textbf{TROPIC01 physical extraction alone is not sufficient, but it is serious.}
If the secure element is physically broken, offline bearer mode should be suspended and authority
rotation should occur unless policy explicitly allows continued risk.

\item \textbf{RP2350 breach can still cause denial of service.}
A breached RP2350 can waste counter steps, corrupt local state, or force online recovery.

\item \textbf{Boot fencing is load bearing for new hardware rejection.}
A copied state image must not be allowed to produce offline releases unless it first proves a boot
chain from the committed boot head.

\item \textbf{Receiver counter evidence is load bearing.}
The receiver must not accept host reported counter state.

\item \textbf{DSM proof verification is load bearing.}
The receiver must verify \(h_i\rightarrow h_{i+1}\). A hardware certificate over an invalid root
does not create value.

\item \textbf{Recovery durability is load bearing for liveness.}
If the counter moved, the same release must remain recoverable. Otherwise the appliance risks an
orphaned commit, which is a liveness failure.

\item \textbf{New relationships remain online checked.}
A closed adversary branch does not affect honest parties because independent relationships and
reconciliation require valid root lineage and counter evidence.

\item \textbf{Tripwire exposes later.}
Tripwire exposes forks on reconciliation.

\item \textbf{Strict frontier continuity requires online re admission after a rejected attempt.}
The physical counter and the appliance root advance when a release is produced, even if the
receiver ultimately rejects that confirm. A rejected attempt therefore opens a gap between the
sender lineage and that particular receiver's adopted frontier, and a later transfer to the same
receiver is refused at the frontier check until an online re admission reconciles the frontier.
This is a liveness limitation, not a safety one: the frontier can only be re admitted online,
under the same authenticated counter evidence that guards genesis adoption, and a first transfer
to any receiver is unaffected because of genesis adoption.
\end{enumerate}

\section{Conclusion}

The boot fenced fused anchor authority is the compact form of the DSM offline bearer appliance.

It does not need a large precommit table. It does not need separate offered, pending, armed, and
released protocol names. The load bearing object is the fused transition:

\[
(h_i,\Ahead_i,\Bhead_b,u_i)
\rightarrow
(h_{i+1},\Ahead_{i+1},\Bhead_{b'},u_i+1).
\]

The previous DSM root commits to the anchor bundle, fused anchor head, boot head, and anchor
counter. The receiver verifies the DSM transition to \(h_{i+1}\). The receiver challenge binds the
release to the actual counterparty. The RP2350 partition certificate proves the partition lineage
participated. TROPIC01 MACANDD proves enrolled hardware witness participation. Authenticated
TROPIC01 counter evidence proves the physical counter corresponds to the exact next anchor counter.
The fused anchor head binds these into a single non interchangeable lineage.

TROPIC01 is necessary but not sufficient. The RP2350 partition is necessary but not sufficient.
DSM validity remains public and receiver verified.

A copied state image cannot resume offline bearer mode on new hardware because offline release
requires a boot ticket chained from the DSM committed boot head under the enrolled anchor bundle.
New hardware requires online authority rotation.

If RP2350 is breached, the attacker can brick the appliance or force online recovery, but cannot
make an honest receiver accept an invalid root transition or a second valid successor from the same
spendable parent without also satisfying DSM verification, boot evidence, partition evidence,
authenticated TROPIC01 counter evidence, and fused anchor head verification.

If the counter moves, the same release must remain recoverable. Recovery therefore re emits the
same committed release and finalization is guarded by the live counter derived anchor counter.

This design is no longer only specified: the complete path --- chip produced release, receiver
operated authenticated counter evidence over a relay, the full acceptance predicate, and canonical
value delivery --- has been exercised end to end on the target hardware.

The whole design reduces to one sentence:

\[
\boxed{
\text{The receiver accepts only a publicly valid DSM root advance whose boot ticket, partition
certificate, TROPIC01 witness, authenticated counter evidence, and fused anchor head all verify
over the same exact successor.}
}
\]

\begin{thebibliography}{9}

\bibitem{dsm}
Cryptskii.
\emph{DSM: Deterministic State Machines}.
Irrefutable Labs Inc., 2026.

\bibitem{tropic}
Tropic Square.
\emph{TROPIC01 Datasheet, API, Application Notes, and Libtropic Documentation}.
Manufacturer documentation.

\bibitem{macandd}
Tropic Square.
\emph{TROPIC01 PIN Verification Application Note}.
ODN TR01 app 002, Version 1.2, 2026.

\bibitem{wots}
Johannes Buchmann, Erik Dahmen, Sarah Ereth, Andreas H\"ulsing, and Markus R\"uckert.
\emph{On the Security of the Winternitz One Time Signature Scheme}.
AFRICACRYPT 2011.

\bibitem{tla}
Leslie Lamport.
\emph{Specifying Systems: The TLA plus Language and Tools for Hardware and Software Engineers}.
Addison Wesley, 2002.

\end{thebibliography}

\end{document}
